%%%%%%%%%%%%%%%%%%%%%%%%%%%%%%%%%%%%%%%%%%%%%%%%%%%%%%%%%%%%%%%%%%%%%%
% Câu chưa có mã dạng — cần gắn %[..] rồi chuyển vào đúng file dạng
% Số câu: 54  (sinh tự động bởi tach_cau.py; từ nay thầy sửa trực tiếp file này)
%%%%%%%%%%%%%%%%%%%%%%%%%%%%%%%%%%%%%%%%%%%%%%%%%%%%%%%%%%%%%%%%%%%%%%

%%%%%%%%%%%%%%%% TRẮC NGHIỆM 4 PHƯƠNG ÁN %%%%%%%%%%%%%%%%

% nguồn: Lop_12_Dai_so_Chuong_6
\begin{ex}%[ma:113829b5]
	 %[2D6H1-2]
	Học sinh lớp 12A tham gia các câu lạc bộ bóng bàn và cờ vua của trường. Chọn ngẫu nhiên 1  học sinh. Xác suất chọn được học sinh tham gia cả hai câu lạc bộ trên bằng $0{,}2$, còn xác suất chọn được học sinh tham gia câu lạc bộ bóng bàn bằng $0{,}25$. Xác suất chọn được học sinh tham gia câu lạc bộ cờ vua, biết học sinh đó tham gia câu lạc bộ bóng bàn, là
	 \choice
		 {\True $0{,}8$}
		{$0{,}45$} 
		 {$0{,}05$}
		{$0{,}2$} 
	\loigiai{ 
		Gọi $A$ là biến cố \lq\lq Chọn được học sinh tham gia câu lạc bộ  cờ vua\rq\rq, $B$ là biến cố \lq\lq Chọn được học sinh tham gia câu lạc bộ bóng bàn\rq\rq.\\
		Ta  có:
		\begin{itemize}
			\item Xác suất chọn được học sinh tham gia cả hai câu lạc bộ: $\mathrm{P}\big(A\cap B\big) =  0{,}2$.
			\item Xác suất chọn được học sinh tham gia câu lạc  bộ bóng bàn: $\mathrm{P}(B) = 0{,}25$.
		\end{itemize}
		Xác suất chọn được học sinh tham gia câu lạc bộ cờ  vua, khi biết học sinh đó tham gia câu lạc bộ bóng bàn là \[\mathrm{P}(A\mid B) = \dfrac{\mathrm{P}(A \cap B)}{\mathrm{P}(B)}= \dfrac{0{,}2}{0{,}25} = 0{,}8.\]
		Vậy, xác suất chọn được học  sinh tham gia câu lạc bộ cờ vua khi đã biết học sinh đó tham gia câu lạc bộ bóng bàn là $0{,}8$.
	} 
\end{ex}

% nguồn: Lop_12_Dai_so_Chuong_6
\begin{ex}%[ma:355ea2d3]
	%[2D6V2-2] 
	Một bộ bài tú lơ khơ gồm $52$ quân bài, trong đó có $4$ quân Át. Bạn Hoa rút ngẫu nhiên $1$ quân bài (không hoàn lại), sau đó bạn Dung rút ngẫu nhiên $1$ quân  bài. Xác suất bạn Dung rút được quân Át là
	 \choice
		{$\dfrac{1}{51}$} 
		 {\True $\dfrac{1}{13}$}
		 {$\dfrac{1}{17}$}
		{$\dfrac{4}{51}$} 
	\loigiai{ 
		 Xét các biến cố:\\
		$A$: \lq\lq Bạn Hoa  rút được quân Át\rq\rq;\\
		$B$: \lq\lq  Bạn Dung rút được quân Át\rq\rq.\\
		Khi  đó, $\mathrm{P}(A)=\dfrac{4}{52}=\dfrac{1}{13}$; $\mathrm{P}\left(\overline{A}\right)=1-\mathrm{P}(A)=1-\dfrac{1}{13}=\dfrac{12}{13}$.
		\begin{itemize}
			\item Nếu bạn Hoa rút được lá Át thì  còn lại $51$ quân bài, trong đó có $3$ lá Át, suy ra $\mathrm{P}(B \mid A)=\dfrac{3}{51}$.
			\item Nếu bạn Hoa không rút được lá Át thì còn lại $51$ quân bài, trong đó có $4$ lá Át, suy ra $\mathrm{P}\left(B \mid \overline{A}\right)=\dfrac{4}{51}$. 
		\end{itemize}
		Theo công thức  xác suất toàn phần, xác suất bạn Dung rút được quân Át là \[\mathrm{P}(B)=\mathrm{P}(A) \cdot \mathrm{P}(B \mid A)+\mathrm{P}\left(\overline{A}\right) \cdot \mathrm{P}\left(B \mid \overline{A}\right)=\dfrac{1}{13} \cdot \dfrac{3}{51}+\dfrac{12}{13} \cdot \dfrac{4}{51}=\dfrac{1}{13}.\]
	} 
\end{ex}

% nguồn: Lop_12_Dai_so_Chuong_6
\begin{ex}%[ma:8dfb6bb2]
	 %[2D6V1-2]
	Một thành phố có $25 \%$ người đàn ông nghiện thuốc lá, trong số những người đàn ông nghiện thuốc lá có $41 \%$ người đàn ông bị bệnh viêm phổi. Chọn ngẫu nhiên một người đàn ông trong thành phố. Xác suất người đàn  ông được chọn bị bệnh viêm phổi, biết người đó nghiện thuốc lá, là
	 \choice
		{$\dfrac{25}{100}$} 
		{$\dfrac{25}{41}$} 
		{$\dfrac{66}{100}$} 
		 {\True $\dfrac{41}{100}$}
	\loigiai{ 
		Gọi $A$ là biến cố \lq\lq Người đàn  ông được chọn bị bệnh bị viêm phổi\rq\rq, $B$ là biến cố \lq\lq Người đàn ông được chọn nghiện thuốc là\rq\rq.\\
		Ta có: 
		\begin{itemize}
			\item Xác suất người đàn ông được chọn vừa nghiện thuốc lá vừa bị viêm phổi: \[\mathrm{P}\big(A\cap B\big) =  0{,}25 \cdot 0{,}41=0{,}1025.\]
			\item Xác suất người đàn ông được chọn nghiện thuốc  lá: $\mathrm{P}(B) = 0{,}25$.
		\end{itemize}
		 Xác suất người đàn ông được chọn bị viêm phổi, biết người đó nghiện thuốc lá \[\mathrm{P}(A\mid B) = \dfrac{\mathrm{P}(A \cap B)}{\mathrm{P}(B)}= \dfrac{0{,}1025}{0{,}25} = 0{,}41.\]
		Vậy xác suất cần  tìm là $\dfrac{41}{100}.$
	 }
\end{ex}

% nguồn: Lop_12_Dai_so_Chuong_6
\begin{ex}%[ma:a827e78a]
	%[2D6V2-2] 
	Có hai chiếc hộp giống nhau. Hộp thứ nhất có $5$ viên bi xanh và $2$ viên bi đỏ. Hộp thứ hai có $4$ viên bi xanh và $6$ viên bi  đỏ. Các viên bi có cùng kích thước và khối lượng. Bạn Sơn chọn ngẫu nhiên một hộp bi và từ đó lấy ra ngẫu nhiên $1$ viên bi. Biết rằng viên bi được lấy ra có màu đỏ, xác suất bạn Sơn chọn hộp thứ hai là
	\choice 
		{$0{,}6$} 
		{$\dfrac{2}{7}$} 
		 {$\dfrac{31}{70}$}
		{\True $\dfrac{21}{31}$} 
	\loigiai{ 
		Gọi $A$ là biến cố viên bi  lấy ra là bi đò.\\
		Gọi  $B$ là biến cố bạn Sơn chọn hộp thứ hai.\\
		Ta có $\mathrm{P}(B)=\mathrm{P}(\overline{B})=0{,}5;  \mathrm{P}(A\mid B)=0{,}6; \mathrm{P}(A \mid \overline{B})=\dfrac{2}{7}$.
		Áp dụng công  thức xác suất toàn phần, ta có
		\[ 
		\mathrm{P}(A)=\mathrm{P}(A  \mid B) \mathrm{P}(B)+\mathrm{P}(A\mid \overline{B}) \mathrm{P}(\overline{B})=0{,}6 \cdot 0{,}5+\dfrac{2}{7} \cdot 0{,}5=\dfrac{31}{70}.
		\] 
		Biết rằng viên bi có màu đỏ, xác suất bạn Sơn chọn  hộp thứ hai là
		\[ 
		\mathrm{P}(B \mid A)=\dfrac{\mathrm{P}(A \mid B) \mathrm{P}(B)}{\mathrm{P}(A)}=\dfrac{0{,}6 \cdot 0{,}5}{\dfrac{31}{70}}=\dfrac{21}{31}.\] 
	} 
\end{ex}

% nguồn: Lop_12_Dai_so_Chuong_6
\begin{ex}%[ma:e015654e]
	 %[2D6V2-2]
	Có hai hộp đựng bi. Hộp thứ nhất có $4$ viên bi xanh và $2$ viên bi đỏ. Hộp thứ hai có $3$ viên bi xanh và $6$ viên bi đỏ. Các viên bi có cùng kích thước và khối lượng. Bạn Thảo chọn ngẫu nhiên một viên bi  từ hộp thứ nhất, bỏ vào hộp thứ hai rồi chọn ra ngẫu nhiên $2$ viên bi từ hộp thứ hai. Xác suất hai viên bi lấy ra từ hộp thứ hai có cùng màu là
	 \choice
		 {$\dfrac{2}{3}$}
		{$\dfrac{8}{15}$} 
		 {$\dfrac{7}{15}$}
		 {\True $\dfrac{22}{45}$}
	 \loigiai{
		Gọi $A_1$ là biến cố viên bi lấy ra từ hộp thứ nhất là bi xanh; $A_2$ là biến cố viên bi lấy ra  từ hộp thứ nhất là bi đỏ.\\
		Gọi $B$ là biến cố hai viên bi được lấy ra từ hộp thứ hai có cùng  màu.\\
		Ta  có
		\[ 
		\mathrm{P}\left(A_{1}\right)=\dfrac{2}{3}; \mathrm{P}\left(A_2\right)=\dfrac{1}{3}; \mathrm{P}\left(B \mid A_{1}\right)=\dfrac{\mathrm{C}_4^2+\mathrm{C}_6^2}{\mathrm{C}_{10}^2}=\dfrac{7}{15};  \mathrm{P}\left(B \mid A_2\right)=\dfrac{\mathrm{C}_3^2+\mathrm{C}_7^2}{\mathrm{C}_{10}^2}=\dfrac{8}{15}.
		 \]
		Áp dụng công thức xác suất toàn phần, ta  có
		\[ 
		\mathrm{P}(B)=\mathrm{P}\left(B \mid A_1\right) \mathrm{P}\left(A_1\right)+\mathrm{P}\left(B \mid  A_2\right) \mathrm{P}\left(A_2\right)=\dfrac{7}{15} \cdot \dfrac{2}{3}+\dfrac{8}{15} \cdot \dfrac{1}{3}=\dfrac{22}{45}.\]
	} 
\end{ex}

% nguồn: Lop_12_Dai_so_Chuong_6
\begin{ex}%[ma:70716384]
	 %[2D6V2-2]
	Khi tìm hiểu về việc học Tiếng Anh của một trường phổ thông, người ta thấy rằng có $70 \%$ học sinh tự học Tiếng Anh bằng hình thức trực tuyến. Chọn ngẫu nhiên $1$ học sinh. Khi đó, xác suất chọn được học sinh giỏi Tiếng Anh, biết học sinh  đó tự học bằng hình thức trực tuyến là $0{,}8$; xác suất chọn được học sinh giỏi Tiếng Anh, biết học sinh đó không tự học bằng hình thức trực tuyến là $0{,}3$. Xác suất chọn được học sinh giỏi Tiếng Anh là
	 \choice
		 {$0{,}24$}
		{$0{,}56$} 
		 {$0{,}7$}
		{\True $0{,}65$} 
	\loigiai{ 
		Gọi $A$ là biến cố \lq\lq Chọn được học sinh giỏi Tiếng Anh\rq\rq, $B$ là biến cố \lq\lq Chọn được học sinh  tự học Tiếng Anh trực tuyến\rq\rq.\\
		Khi  đó, $\mathrm{P}(B)=0{,}7$; $\mathrm{P}\left(\overline{B}\right)=1-0{,}7=0{,}3$.\\
		Xác suất học sinh giỏi Tiếng Anh, biết học sinh đó tự học bằng trực tuyến: $\mathrm{P}(A\mid B) =  0{,}8$.\\
		Xác suất học sinh  giỏi Tiếng Anh, biết học sinh đó không tự học bằng trực tuyến: $\mathrm{P}\left(A\mid \overline{B}\right) = 0{,}3$.\\
		Theo công thức xác suất toàn phần, xác suất chọn được học  sinh giỏi Tiếng Anh là \[\mathrm{P}(A)=\mathrm{P}(B)\cdot \mathrm{P}(A\mid B)+\mathrm{P}\left(\overline{B}\right)\cdot \mathrm{P}\left(A\mid \overline{B}\right)=0{,}7\cdot 0{,}8 + 0{,}3\cdot 0{,}3=0{,}65.\]
	} 
\end{ex}

% nguồn: Lop_12_Dai_so_Chuong_6
\begin{ex}%[ma:cbf966e3]
	 %[2D6V2-2]
	Có hai chiếc hộp giống nhau. Hộp thứ nhất có $5$ viên bi xanh và $2$ viên bi đỏ. Hộp thứ hai có $4$ viên bi xanh và $6$ viên bi đỏ. Các viên bi có cùng kích thước và khối lượng. Bạn Sơn chọn ngẫu nhiên một hộp  bi và từ đó lấy ra ngẫu nhiên $1$ viên bi. Xác suất viên bi được chọn có màu xanh là
	 \choice
		{$\dfrac{9}{17}$} 
		{$\dfrac{2}{7}$} 
		 {$\dfrac{39}{140}$}
		{\True  $\dfrac{39}{70}$}
	 \loigiai{
		Gọi $A_1$ là  biến cố bạn Sơn chọn hộp thứ nhất; $A_2$ là biến cố bạn Sơn chọn hộp thứ hai.\\
		Gọi $B$ là biến cố  viên bi được chọn có màu xanh.\\
		Ta có  $\mathrm{P}\left(A_{1}\right)=\mathrm{P}\left(A_{2}\right)=0{,}5; \mathrm{P}\left(B \mid A_{1}\right)=\dfrac{5}{7}; \mathrm{P}\left(B \mid A_{2}\right)=\dfrac{4}{10}$.\\
		 Áp dụng công thức xác suất toàn phần, ta có
		 \[\mathrm{P}(B)=\mathrm{P}\left(B \mid A_{1}\right) \mathrm{P}\left(A_{1}\right)+\mathrm{P}\left(B \mid A_{2}\right) \mathrm{P}\left(A_{2}\right)=\left(\dfrac{5}{7}+\dfrac{4}{10}\right) \cdot \dfrac{1}{2}=\dfrac{39}{70}.\]
	 }
\end{ex}

% nguồn: Lop_12_Dai_so_Chuong_6
\begin{ex}%[ma:af4c0fde]
	%[2D6H1-2] 
	Khi  điều tra về hoạt động sử dụng máy tính và tình trạng cận thị của trẻ em ở một tỉnh thì được kết quả:
	\begin{enumerate}[-]
		 \item Có $10\%$ trẻ em thường xuyên sử dụng máy tính;
		\item Có $30\%$  trẻ em bị cận thị;
		\item Trong  những trẻ em thường xuyên sử dụng máy tính có $54\% $ trẻ em bị cận thị.
	\end{enumerate}
	Chọn ngẫu nhiên $1$ trẻ em. Xác suất trẻ em được chọn thường  xuyên sử dụng máy tính, biết trẻ em đó bị cận thị, là
	 \choice
		 {$0{,}94$}
		 {$0{,}14$}
		 {\True$0{,}18$}
		 {$0{,}0162$}
	\loigiai{ 
		Gọi $A$ là biến cố \lq\lq Trẻ em thường xuyên sử dụng máy tính\rq\rq.  \\
		Ta có  $\mathrm{P}(A)=0{,}1 >0$.\\
		Gọi $B$ là biến cố \lq\lq Trẻ em bị  cận thị\rq\rq.\\
		Ta  có $\mathrm{P}(B)=0{,}3>0$.\\
		Xác suất trẻ em bị cận thị, biết thường xuyên sử dụng máy tính là  $\mathrm{P}(B \mid A)=0{,}54$.\\
		Ta có $\mathrm{P}(B  \mid A)=\dfrac{\mathrm{P}(BA)}{\mathrm{P}(A)} \Rightarrow \mathrm{P}(BA)=0{,}54 \cdot 0{,}1=0{,}054$.\\
		 Do đó: $\mathrm{P}(A \mid B)=\dfrac{\mathrm{P}(BA)}{\mathrm{P}(B)}=\dfrac{0{,}054}{0{,}3}=0{,}18$.
	} 
\end{ex}

% nguồn: Lop_12_Dai_so_Chuong_6
\begin{ex}%[ma:15612625]
	 %[2D6H1-2]
	Một thư viện có hai phòng riêng biệt, phòng A và phòng B. Xác suất chọn được một quyển sách về chủ đề Khoa học tự nhiên thuộc phòng A và thuộc phòng B lần lượt là $0{,}25$ và  $0{,}5$. Chọn ngẫu nhiên $1$ quyển sách của thư viện. Giả sử quyển sách được chọn về chủ đề Khoa học tự nhiên, xác suất quyển sách đó ở phòng A là
	\choice 
		{$\dfrac{2}{3}$} 
		{$\dfrac{1}{2}$} 
		{\True  $\dfrac{1}{4}$}
		{$\dfrac{1}{3}$} 
	 \loigiai{
		Xét các biến cố 
		\begin{itemize}
			\item $M\colon$\lq\lq Quyển sách được chọn ở phòng  A\rq\rq;
			\item $N\colon$\lq\lq Quyển sách  được chọn về chủ đề Khoa học tự nhiên\rq\rq;
			 \item $Q\colon$\lq\lq Quyển sách được chọn về chủ đề Khoa học tự nhiên và thuộc phòng A\rq\rq;
			\item $R\colon$\lq\lq Quyển sách được chọn về chủ đề Khoa học tự nhiên và thuộc  phòng B\rq\rq.
		\end{itemize}
		Nhận thấy $N=Q \cup R$  và $Q$, $R$ là hai biến cố xung khắc nên
		\[\mathrm{P}(N)=\mathrm{P}(Q)+\mathrm{P}(R)=0{,}25+0{,}5=0{,}75.\] 
		 Ta có $\mathrm{P}(M \mid N)=\dfrac{\mathrm{P}(M \cap N)}{\mathrm{P}(N)}=\dfrac{0{,}25}{0{,}75}=\dfrac{1}{3}$.\\
		Vậy xác suất quyển sách được chọn ở phòng $\mathrm{A}$, biết rằng quyển sách đó về chủ đề Khoa học  tự nhiên, là $\dfrac{1}{3}$.
	 }
\end{ex}

% nguồn: Lop_12_Dai_so_Chuong_6
\begin{ex}%[ma:a14cf143]
	 %[2D6H1-2]
	Một động cơ điện có hai van bảo hiểm cùng hoạt động. Xác suất hoạt động tốt của van I là $0{,}9$, của van II là $0{,}72$. Xác suất hoạt động tốt của van I, biết van II hoạt động tốt là $0{,}96$. Giả sử van I hoạt động tốt, xác suất  hoạt động tốt của van II là
	 \choice
		 {$0{,}675$}
		{\True $0{,}768$} 
		{$0{,}66$} 
		{$0{,}78$} 
	\loigiai{ 
		Gọi $A$ là  biến cố \lq\lq Van I hoạt động tốt\rq\rq. Ta có $ \mathrm{P}(A)=0{,}9$.\\
		Gọi $B$ là biến cố \lq\lq van II  hoạt động tốt\rq\rq. Ta có $\mathrm{P}(B)=0{,}72$.\\
		Gọi $C$ là biến cố \lq\lq Van I hoạt động tốt biết van II hoạt động tốt\rq\rq. Ta  có $ \mathrm{P}(C)=0{,}96$.\\
		 Ta có $\mathrm{P}(B)>0$
		nên \[\mathrm{P}(C)=\mathrm{P}(A \mid B)=\dfrac{\mathrm{P}(AB)}{\mathrm{P}(B)} \Rightarrow \mathrm{P}(AB)=\mathrm{P}(A \mid B)\cdot \mathrm{P}(B)=0{,}96  \cdot 0{,}72=0{,}6912.\]
		Vậy giả sử van I hoạt động tốt, xác suất  van II hoạt động tốt là\[\mathrm{P}(B \mid A)=\dfrac{\mathrm{P}(BA)}{\mathrm{P}(A)}=\dfrac{0{,}6912}{0{,}9}=0{,}768.\]
	 }
\end{ex}


%%%%%%%%%%%%%%%% ĐÚNG/SAI %%%%%%%%%%%%%%%%

% nguồn: Lop_12_Dai_so_Chuong_6
\begin{ex}%[ma:ae25b014]
	 %[2D6H2-3]
	Một kho hàng có $85 \%$ sản phẩm loại I và $15 \%$ sản phẩm loại II, trong đó có $1 \%$ sản phẩm loại I bị hỏng, $4 \%$  sản phẩm loại II bị hỏng. Các sản phẩm có kích thước và hình dạng như nhau. Một khách hàng chọn ngẫu nhiên $1$ sản phẩm. Xét các biến cố:\\
	$A\colon$\lq\lq Khách  hàng chọn được sản phẩm loại I\rq\rq;\\
	$B\colon$\lq\lq  Khách hàng chọn được sản phẩm không bị hỏng\rq\rq.
	 \choiceTF
		 {\True $\mathrm{P}(A)=0{,}85$}
		{\True  $\mathrm{P}(B \mid A)=0{,}99$}
		{\True  $\mathrm{P}(B)=0{,}9855$}
		{$\mathrm{P}(A \mid  B)=0{,}95$}
	\loigiai{ 
		\begin{itemchoice}
			 \itemch {\bf Đúng}.\\ Vì kho hàng có $85 \%$ sản phẩm loại I nên $\mathrm{P}(A)=0{,}85$.
			\itemch {\bf Đúng}.\\ Vì kho  hàng có $1 \%$ sản phẩm loại I bị hỏng nên $\mathrm{P}(B \mid A)=1-\mathrm{P}\big(\overline{B}\mid A\big)=1-0{,}01=0{,}99$.
			\itemch {\bf  Đúng}.\\ Vì $\mathrm{P}(B)=\mathrm{P}(A)\cdot \mathrm{P}(B\mid A)+\mathrm{P}\big(\overline{A}\big)\cdot \mathrm{P}(B\mid \overline{A})=0{,}85\cdot 0{,}99+0{,}15\cdot 0{,}96=0{,}9855$.
			\itemch {\bf Sai}.\\ 
			Theo công thức  Bayes ta có
			\[\mathrm{P}(A \mid B)=\dfrac{\mathrm{P}\big(A\big)\cdot \mathrm{P}\big(B\mid  A\big)}{\mathrm{P}(B)}=\dfrac{0{,}85\cdot0{,}99}{0{,}9855}\approx0{,}854.\]
		\end{itemchoice}
	 }
\end{ex}

% nguồn: Lop_12_Dai_so_Chuong_6
\begin{ex}%[ma:fed01ae3]
	 %[2D6H2-3]
	Một xưởng máy sử dụng một loại linh kiện được sản xuất từ hai cơ sở I và II. Số linh kiện do cơ sở I sản xuất chiếm $61 \%$, số linh kiện do cơ sở II sản xuất chiếm $39 \%$. Tỉ lệ linh kiện đạt tiêu chuẩn của cơ sở I, cơ sở II lần lượt là $93 \%$, $82 \%$. Kiểm tra ngẫu nhiên $1$ linh kiện ở  xường máy. Xét các biến cố:\\
	$A_1\colon$\lq\lq Linh  kiện được kiểm tra do cơ sở I sản xuất\rq\rq;\\
	 $A_2\colon$\lq\lq Linh kiện được kiểm tra do cơ sở II sản xuất\rq\rq;\\
	$B\colon$\lq\lq Linh  kiện được kiểm tra đạt tiêu chuẩn\rq\rq.
	\choiceTF 
		 {$\mathrm{P}\left(A_{1}\right)=0{,}39$}
		 {\True $\mathrm{P}\left(B \mid A_{2}\right)=0{,}82$}
		{\True  $\mathrm{P}(B)=0{,}8871$}
		 {$\mathrm{P}\left(A_{1} \mid B\right)=0{,}55$}
	 \loigiai{
		\begin{itemchoice}
			\itemch {\bf Sai}.\\ Vì số linh kiện do cơ sở I sản  xuất chiếm $61 \%$ nên $\mathrm{P}\left(A_{1}\right)=0{,}65$.
			\itemch {\bf  Đúng}.\\ Vì tỉ lệ linh kiện đạt tiêu chuẩn của cơ sở II là $82 \%$ nên $\mathrm{P}\left(B \mid A_{2}\right)=0{,}82$.
			\itemch {\bf Đúng}.\\ Theo công thức xác suất toàn phần ta có  \[\mathrm{P}(B)=\mathrm{P}(A_1)\cdot \mathrm{P}(B\mid A_1)+\mathrm{P}(A_2)\cdot \mathrm{P}(B\mid A_2)=0{,}61\cdot 0{,}93+0{,}39\cdot 0{,}82=0{,}8871.\]
			\itemch  {\bf Sai}.\\
			Theo công thức  Bayes ta có
			\[\mathrm{P}\big(A_1\mid B\big)=\dfrac{\mathrm{P}\big(A_1\big)\cdot \mathrm{P}\big(B\mid  A_1\big)}{\mathrm{P}(B)}=\dfrac{0{,}65\cdot 0{,}93}{0{,}8871}\approx0{,}64.\]
		\end{itemchoice}
	 }
\end{ex}

% nguồn: Lop_12_Dai_so_Chuong_6
\begin{ex}%[ma:f6f5794a]
	%[2D6H1-2] 
	Khi kiểm tra sức khoẻ tổng quát của bệnh nhân  ở một bệnh viện, người ta được kết quả như sau
	\begin{itemize}
		\item Có $40\%$ bệnh nhân bị đau  dạ dày
		\item Có $30\%$ bệnh  nhân thường xuyên bị stress
		\item Trong số các bệnh nhân bị stress có $80\%$ bệnh nhân  bị đau dạ dày
	\end{itemize}
	Chọn ngẫu  nhiên $1$ bệnh nhân.
	\choiceTF 
		{\True Xác suất chọn được bệnh nhân thường xuyên  bị stress là~$0{,}3$}
		{\True Xác suất chọn được bệnh nhân bị đau dạ dày, biết bệnh nhân đó thường xuyên bị stress  là~$0{,}8$}
		{\True Xác suất chọn được bệnh nhân vừa thường xuyên bị  stress vừa bị đau dạ dày là~$0{,}24$}
		{\True Xác  suất chọn được bệnh nhân thường xuyên bị stress, biết bệnh nhân đó bị đau dạ dày là~$0{,}6$}
	\loigiai{\begin{itemchoice}
	 \itemch
	Xét biến cố: $A$: \lq\lq Chọn được bệnh nhân thường  xuyên bị stress\rq\rq. Ta có $\mathrm{P}(A) = 0{,}3$.
	\itemch 
	Biến cố $B$: \lq\lq Chọn được  bệnh nhân bị đau dạ dày\rq\rq, $\mathrm{P}(B) = 0{,}4$.\\
	Biến cố $B\mid A$: \lq\lq  Bệnh nhân bị stress có bị đau dạ dày\rq\rq. Ta có $\mathrm{P}(B\mid A) = 0{,}8$.
	\itemch 
	Xác suất chọn được bệnh nhân thường xuyên bị stress vừa bị đau dạ dày là~$\mathrm{P}\left(A \cap B\right) = \mathrm{P}(A)\cdot  \mathrm{P}(B\mid A)=0{,}3 \cdot 0{,}8 = 0{,}24$.
	 \itemch
	 Xác suất chọn được bệnh nhân thường xuyên bị stress, biết bệnh nhân đó bị đau dạ dày là
	$\mathrm{P}\left(A\mid B\right) =\dfrac{\mathrm{P}\left(A\cap B\right)}{\mathrm{P}(B)}= \dfrac{0{,}24}{0{,}4} = 0{,}6$. 
\end{itemchoice}
} 
\end{ex}

% nguồn: Lop_12_Dai_so_Chuong_6
\begin{ex}%[ma:0a472f4b]
	 %[2D6H2-2]
	Trong một hộp có $10$ quả bóng màu xanh và $12$ quả bóng màu đỏ, các quả bóng có khối lượng và kích thước như nhau. Bạn Tuấn lấy ngẫu nhiên lần lượt $2$ quả bóng, mỗi lần lấy $1$ quả và không hoàn  lại. Xét các biến cố:\\
	 $A\colon$\lq\lq Lần thứ nhất lấy được quả bóng màu xanh\rq\rq;\\
	$B\colon$\lq\lq Lần  thứ hai lấy được quả bóng màu xanh\rq\rq.
	\choiceTF 
		{\True  $\mathrm{P}(A)=\dfrac{5}{11}$}
		{$\mathrm{P}(B \mid  A)=\dfrac{10}{21}$}
		 {$\mathrm{P}\big(B \mid \overline{A}\big)=\dfrac{3}{7}$}
		{\True $\mathrm{P}(B)=\dfrac{5}{11}$} 
	 \loigiai{
		\begin{itemchoice}
			\itemch {\bf Đúng}.\\ 
			Số kết quả thuận lợi cho biến cố $A$  là $n(A)=10\cdot 9 +10\cdot 12$.\\
			Không  gian mẫu là $n(\Omega)=22\cdot 21$.\\
			Vậy $\mathrm{P}(A)=\dfrac{n(A)}{n(\Omega)}=\dfrac{10  \cdot 9+10\cdot 12}{22 \cdot 21}=\dfrac{5}{11}$.
			\itemch  {\bf Sai}.\\
			Ta  có biến cố $A\cap B\colon$\lq\lq Cả hai lần đều lấy được quả bóng màu xanh\rq\rq.\\
			Số kết quả thuận lợi cho biến  cố $A\cap B$ là $n(A\cap B)=10\cdot 9$.\\
			Vậy $\mathrm{P}(B \mid A)=\dfrac{\mathrm{P}\big(A\cap B\big)}{\mathrm{P}(A)}=\dfrac{n(A\cap  B)}{n(A)}=\dfrac{10\cdot 9}{10\cdot 9+10\cdot 12}=\dfrac{3}{7}$.
			\itemch  {\bf Sai}.\\
			Ta có biến cố $\overline{A}\cap B\colon$\lq\lq Lần thứ nhất lấy được quả bóng màu đỏ  và lần thứ hai lấy được quả bóng màu xanh\rq\rq.\\
			Số kết quả thuận lợi cho biến cố $\overline{A}\cap B$  là $n\big(\overline{A}\cap B\big)=12\cdot 10$.\\
			Số kết quả thuận lợi cho biến cố $\overline{A}$ là $n\big(\overline{A}\big)=12\cdot 10+12\cdot 11$ (ứng với hai trường hợp  đỏ - xanh và đỏ - đỏ).\\
			Vậy $\mathrm{P}\big(B \mid \overline{A}\big)=\dfrac{\mathrm{P}(\overline{A} \cap B)}{\mathrm{P}\big(\overline{A}\big)}=\dfrac{n\big(\overline{A}\cap B\big)}{n\big(\overline{A}\big)}=\dfrac{12 \cdot 10}{12 \cdot 10+12\cdot 11}=\dfrac{10}{21}$. 
			\itemch {\bf Đúng}.\\ 
			Số kết quả thuận lợi cho biến  cố $B$ là $n(B)=12\cdot 10+10\cdot 9$ (ứng với hai trường hợp đỏ - xanh và xanh - xanh).\\
			Vậy $\mathrm{P}(B)=\dfrac{12 \cdot  10+10\cdot 9}{22 \cdot 21}=\dfrac{5}{11}$.
		\end{itemchoice}
	 }
\end{ex}

% nguồn: Lop_12_Dai_so_Chuong_6
\begin{ex}%[ma:fa3adede]
	 %[2D6C1-4]
	Sau học kì I năm học 2024-2025, thầy Nghĩa chủ nhiệm lớp 12B5 nhận thấy rằng lớp mình có $60\%$ học sinh có kết quả xuất sắc, $40\%$ học sinh có kết quả loại giỏi,  không có học sinh khá và trung bình. Nhưng để nắm chính xác hơn về năng lực tư duy môn toán của từng học sinh nên thầy Nghĩa đã cho học sinh làm bài kiểm tra toán trong $90$ phút. Sau khi chấm bài xong, thầy Nghĩa thấy rằng trong số học sinh loại giỏi có $8$ học sinh từ $9$ điểm toán trở lên và có $75\%$ học sinh xuất sắc trong các học sinh được điểm toán từ $9$ trở lên. Biết lớp 12B5 có $40$ học sinh.
	\choiceTF 
		{\True Tỉ lệ học sinh có điểm toán từ $9$ trở lên  của lớp 12B5 là~$80\%$}
		{\True Học sinh xuất sắc kiểm tra môn toán  đều lớn hơn hoặc bằng $9$ điểm}
		{\True Những học sinh có điểm toán dưới $9$ điểm đều  là~học sinh loại giỏi}
		{\True Có $22$ học sinh kết quả xuất sắc có điểm trên $9$ biết rằng tỉ lệ học sinh có điểm toán trên $9$ điểm của học sinh giỏi bằng $37,5\%$ và trong số học sinh có điểm  bằng $9$ có $50\%$ học sinh xuất sắc}
	 \loigiai{
		\begin{itemize}
			\item Tổng số  học sinh lớp 12B5 là~$40$.
			\item Số học sinh xuất sắc $60\%$ của $40$  là~$0{,}6 \cdot 40 = 24$.
			\item Số học  sinh loại giỏi $40\%$ của $40$ là~$0{,}4 \cdot 40 = 16$.
		\end{itemize}
		\begin{itemize}
			\item Trong nhóm học sinh loại giỏi, có $8$ học  sinh đạt điểm từ $9$ trở lên.
			\item Trong các học  sinh đạt điểm từ $9$ trở lên, có $75\%$ là~học sinh xuất sắc.
		\end{itemize}
		Gọi $x$ là~số học sinh có điểm từ $9$ trở lên. Do có $8$ học sinh loại giỏi trong nhóm này, nên số học sinh xuất  sắc từ $9$ trở lên là~$(x - 8)$. Theo đề bài, tỉ lệ học sinh xuất sắc trong nhóm này là~$75\%$, ta lập phương trình $\dfrac{x - 8}{x} \;=\; 0{,}75\Leftrightarrow x=32$.\\
		Vậy có $32$ học sinh đạt điểm từ $9$ trở lên. Số học  sinh có điểm dưới $9$ sẽ là~$40-32=8$.\\
		\begin{itemize}
			\item  \textit{Học sinh xuất sắc}
			Tổng số $24$. Số học sinh xuất sắc trong nhóm $\geq 9$ đúng  bằng $32 - 8 = 24$. Như vậy, toàn bộ $24$ học sinh xuất sắc đều đạt $\geq 9$ điểm.
			 \item \textit{Học sinh loại giỏi}
			Tổng số $16$.  Đã biết trong đó có $8$ học sinh đạt $\ge 9$, nên còn $16 - 8 = 8$ học sinh loại giỏi đạt $<9$.
		\end{itemize}
		\begin{itemchoice}
			\itemch \textbf{Đúng} \textit{Tỉ lệ học sinh có điểm toán  từ $9$ trở lên là~$80\%$.}\\
			Ta có tổng số học sinh  là~$40$, trong đó $32$ học sinh $\ge 9$, nên tỉ lệ $\dfrac{32}{40} \;=\; 0{,}8 \;=\; 80\%$.
			\itemch \textbf{Đúng} \textit{Học sinh xuất sắc kiểm tra môn toán  đều $\ge 9$ điểm.}\\
			Vì tất  cả $24$ học sinh xuất sắc nằm trong nhóm $32$ học sinh $\ge 9$, không ai dưới $9$.
			\itemch \textbf{Đúng} \textit{Những học sinh có điểm toán dưới $9$ điểm đều là~học  sinh loại giỏi.}\\
			Có $8$ học sinh dưới $9$ và tất cả họ đều thuộc nhóm loại  giỏi (vì toàn bộ xuất sắc đều $\ge 9$).
			\itemch \textbf{Đúng} \textit{Có $22$ học sinh xuất sắc có điểm $>9$, với các dữ kiện bổ  sung:}
			\begin{itemize}
				\item  Tỉ lệ học sinh loại giỏi có điểm $>9$ là~$37,5\%$. Vậy số học sinh loại giỏi có điểm $>9$ là~$0{,}375 \cdot 16 = 6$.
				\item Do nhóm học sinh loại giỏi có $\ge 9$ là~$8$, trong đó $6$ người  $>9$ nên còn $2$ người đúng $9$ điểm.
				\item  Trong nhóm học sinh xuất sắc $\ge 9$ (tổng $24$ người), giả sử có $x$ người $>9$, vậy còn $24 - x$ người đúng $9$.
				\item Gọi $b$ là~tổng số học sinh có điểm $=9$. Khi đó $b$ gồm $2$ giỏi  và $(24 - x)$ xuất sắc, nghĩa là~$b =2+(24 - x)=26 - x$.
				\item Ngoài ra, đề cho  biết trong số học sinh có điểm $=9$, có $50\%$ là~xuất sắc, nên:
				$$\dfrac{24 - x}{26  - x}= 0{,}5\Leftrightarrow x=22.$$
				Vậy có $22$ học sinh  xuất sắc đạt điểm lớn hơn $9$.
			\end{itemize}
		\end{itemchoice}
	 }
\end{ex}

% nguồn: Lop_12_Dai_so_Chuong_6
\begin{ex}%[ma:cf4c8293]
	%[2D6V2-4] 
	Xác suất để công ty $X$ thuê một trong hai công ty vệ sinh $A$ và $B$ tư vấn lần lượt là~$0{,}4$ và $0{,}6$. Theo kinh nghiệm khả năng $X$ phát sinh thêm chi phí khi  sử dụng dịch vụ tư vấn của công ty $A$ và $B$ lần lượt là~$0{,}05$ và $0{,}03$.
	 \choiceTF
		{\True Xác suất để $X$ có phát sinh thêm chi phí khi sử dụng dịch vụ tư  vấn là~$0{,}038$}
		{Biết $X$ có phát sinh thêm chi  phí khi sử dụng dịch vụ tư vấn. Xác suất để $X$ thuê công ty $A$ tư vấn là~$0{,}4737$}
		{Biết $X$ có phát sinh thêm chi phí khi sử dụng dịch vụ tư vấn. Xác suất để $X$ thuê công ty  $B$ tư vấn là~$0{,}5263$}
		{\True Biết $X$  không phát sinh thêm chi phí khi sử dụng dịch vụ tư vấn. Xác suất để $X$ thuê công ty $A$ tư vấn là~$0{,}395$}
	\loigiai{\begin{itemchoice}
	\itemch \textbf{Đúng}.\\ 
	Gọi $T$ là~biến cố  \lq\lq công ty $X$ thuê công ty $A$\rq\rq.\\
	Khi đó $\overline{T}$ là~biến cố \lq\lq công ty $X$ thuê công ty $B$\rq\rq.\\ 
	Gọi $C$ là~biến cố \lq\lq công ty $X$ phát sinh thêm  chi phí\rq\rq.\\
	Xác suất để $X$ có phát  sinh thêm chi phí là
	\begin{align*}
		\mathrm{P}(C) &= \mathrm{P}(T) \cdot \mathrm{P}\left(C \mid T \right) + \mathrm{P}\left(\overline{T} \right) \cdot  \mathrm{P}\left(C \mid \overline{T} \right)\\
		 &= 0{,}4 \cdot 0{,}05 + 0{,}6 \cdot 0{,}03 = 0{,}038.
	\end{align*}
	\itemch \textbf{Sai}.\\ 
	Ta có $\mathrm{P}\left(T \mid C \right) = \dfrac{\mathrm{P}(T) \cdot \mathrm{P}\left(C \mid T  \right)}{\mathrm{P}(C)} = \dfrac{0{,}05 \cdot 0{,}4}{0{,}038} = 0{,}5263$.
	 \itemch \textbf{Sai}.\\
	Ta có $\mathrm{P}\left(\overline{T} \mid C \right) = 1 - \mathrm{P}\left(T \mid  C \right) = 1-0{,}5263 = 0{,}4737$.
	\itemch \textbf{Đúng}.\\ 
	Ta có $\mathrm{P}(C) = 0{,}038$ suy ra $\mathrm{P}\left(\overline{C}  \right) = 0{,}962$.\\
	Khi đó $\mathrm{P}\left(T \mid \overline{C} \right) = \dfrac{\mathrm{P}\left(\overline{C} \mid T \right)  \cdot \mathrm{P}(T)}{\mathrm{P}\left(\overline{C} \right)} = \dfrac{\left(1-0{,}05 \right) \cdot 0{,}4}{0{,}962} = 0{,}395$.
\end{itemchoice}
} 
\end{ex}

% nguồn: Lop_12_Dai_so_Chuong_6
\begin{ex}%[ma:9a68b0ac]
	 %[2D6H1-2]
	Gieo  một xúc xắc cân đối và đồng chất $1$ lần. Xét các biến cố:\\
	$A\colon$\lq\lq Mặt xuất hiện của  xúc xắc ghi số $5$\rq\rq;\\
	$B\colon$\lq\lq Mặt xuất hiện của xúc xắc  ghi số lẻ\rq\rq.
	 \choiceTF
		{$\mathrm{P}(A)=\dfrac{5}{6}$} 
		 {\True $\mathrm{P}(A \cap B)=\dfrac{1}{6}$}
		 {\True $\mathrm{P}(B \mid A)=1$}
		 {$\mathrm{P}(A \mid B)=\dfrac{1}{2}$}
	 \loigiai{
		\begin{itemchoice}
			\itemch  {\bf Sai}.\\
			Không gian mẫu là  $6$.\\
			 Số kết quả thuận lợi cho biến cố $A$ là $1$.\\
			Vậy $\mathrm{P}(A)=\dfrac{1}{6}$. 
			 \itemch {\bf Đúng}.\\
			 Số kết quả thuận lợi cho biến cố $B$ là $3$.\\
			Số kết quả thuận lợi cho biến cố $A\cap B$ là $1$ (ứng với mặt xuất hiện ghi  số $5$).\\
			Vậy $\mathrm{P}(A  \cap B)=\dfrac{1}{6}$.
			\itemch {\bf Đúng}.\\ Ta có $\mathrm{P}(B  \mid A)=\dfrac{\mathrm{P}(A \cap B)}{\mathrm{P}(A)}=1$.
			\itemch {\bf Sai}.\\ 
			Xác suất xảy ra biến cố $B$ là  $\mathrm{P}\big(B\big)=\dfrac{3}{6}=\dfrac{1}{2}$.\\
			Vậy $\mathrm{P}(A \mid  B)=\dfrac{\mathrm{P}(A \cap B)}{\mathrm{P}(B)}=\dfrac{\dfrac{1}{6}}{\dfrac{1}{2}}=\dfrac{1}{3}$.
		\end{itemchoice}
	 }
\end{ex}

% nguồn: Lop_12_Dai_so_Chuong_6
\begin{ex}%[ma:5d64142e]
	 %[2D6H2-4]
	Một xưởng máy sử dụng một loại linh kiện được sản xuất từ hai cơ sở I và II. Số linh kiện do cơ sở I sản xuất chiếm $61$\%, số linh kiện do cơ sở II sản xuất chiếm  $39$\%. Tỉ lệ linh kiện đạt tiêu chuẩn của cơ sở I, cơ sở II lần lượt là~$93$\%, $82$\%. Kiểm tra ngẫu nhiên một linh kiện ở xưởng máy. Xét các biến cố\\
	$A_1\colon$ \lq\lq Linh kiện  được kiểm tra do cơ sở I sản xuất\rq\rq;\\
	$A_2\colon$ \lq\lq Linh kiện được kiểm tra do cơ sở II  sản xuất\rq\rq;\\
	$B\colon$ \lq\lq Linh kiện  được kiểm tra đạt tiêu chuẩn\rq\rq .
	\choiceTF 
		 {\True Xác suất $\mathrm{P}\left(A_1 \right)=0{,}61$}
		{\True Xác suất có điều kiện $\mathrm{P}\left(B\mid A_2  \right)=0{,}82$}
		{\True Xác suất  $\mathrm{P}\left(B\right)=0{,}8871$}
		 {Xác suất có điều kiện $\mathrm{P}\left(A_1\mid B\right)=0{,}55$}
	 \loigiai{
		\begin{itemchoice}
			\itemch  Do $\mathrm{P}\left(A_1 \right)=0{,}61$.
			\itemch $\mathrm{P}\left(BA_2 \right)=\dfrac{P\left(B\cap A_2 \right)}{P\left(A_2 \right)}=0{,}82$. Do  đó.
			\itemch  Ta có $\mathrm{P}\left(A_1 \right)=0{,}61$; $P\left(A_2 \right)=0{,}39; P\left(BA_1 \right)=0{,}93$; $P\left(BA_2 \right)=0{,}82$.
			Theo công thức xác suất toàn phần, ta có 
			\centerline {$P\left(B\right)=P\left(A_1\right)\cdot P\left(BA_1\right)+P\left(A_2\right)\cdot P\left(BA_2\right)=0{,}61\cdot 0{,}93+0{,}39\cdot  0{,}82=0{,}8871$.}
			\itemch Theo công thức Bayes, ta có $\mathrm{P}\left(A_1  B\right)=\dfrac{P\left(A_1 \right)\cdot P\left(BA_1 \right)}{P\left(B\right)}=\dfrac{0{,}61\cdot 0{,}93}{0{,}8871} \approx 0{,}64$.
		\end{itemchoice}
	 }
\end{ex}

% nguồn: Lop_12_Dai_so_Chuong_6
\begin{ex}%[ma:36e9b123]
	%[2D6H2-1] 
	Năm $2020$, dịch COVID-19 bùng phát trên toàn thế giới. Các nhà khoa  học đã phát triển một loại test nhanh để phát hiện virus SARS-CoV-2 gây bệnh COVID-19. Theo thống kê, khi một người nhiễm virus SARS-CoV-2 thì xác suất để test nhanh cho kết quả dương tính là~$90 \%$. Tuy nhiên, khi một người không bị nhiễm virus, xác suất để test nhanh vẫn cho kết quả dương tính là~$5 \%$. Biết rằng tỷ lệ người nhiễm virus SARV-CoV-2 ở một quốc gia là~$2 \%$ trong dân số.\\
	Gọi $X$ là~biến cố \lq \lq Một người nhiễm virus SARS-CoV-2 \rq \rq \ và  Y là~biến cố \lq \lq một người có kết quả test nhanh dương tính \rq \rq.
	\choiceTF 
		 {\True $\mathrm{P}(X)=0{,}02$}
		{\True $\mathrm{P}(Y\mid X)=0{,}9$} 
		 {\True $\mathrm{P}(X\mid Y)=0{,}567$}
		{$\mathrm{P}(X  \cap Y)=0{,}06$}
	\loigiai{ 
		\begin{itemchoice}
			\itemch Đúng. Ta có  $\mathrm{P}(X)$ là~xác suất một người nhiểm virus SARS-CoV-2. Theo đề bài, tỷ lệ người nhiểm virus ở một quốc gia là~$2\%=0{,}02$ trong dân số.
			\itemch Đúng. $\mathrm{P}(Y\mid X)$ là~xác suất một người có kết quả test nhanh dương tính, với điều kiện người đó nhiễm virus. Theo giả thiết, khi có một người nhiễm  virus, xác suất để test nhanh có kết quả dương tính là~$90 \%=0{,}9$.
			\itemch Đúng.  Ta có $\mathrm{P}(Y \mid X)=0{,}9$, $\mathrm{P}(X)=0{,}02$.\\
			Suy ra  $\mathrm{P}(Y)=\mathrm{P}(A)\cdot \mathrm{P}(Y \mid A)+\mathrm{P}(\overline{X})\cdot \mathrm{P}(Y \mid \overline{X})=0{,}9 \cdot 0{,}02+0{,}98 \cdot 0{,}05=0{,}0634$.
			\itemch  Sai. Ta có $\mathrm{P}(Y\cap X)=\mathrm{P}(X)\cdot \mathrm{P}(Y \mid X)=0{,}02 \cdot 0{,}9=0{,}05$.
		\end{itemchoice}
	} 
\end{ex}

% nguồn: Lop_12_Dai_so_Chuong_6
\begin{ex}%[ma:f0178816]
	 %[2D6V2-2]
	Một loại xét nghiệm nhanh SARS-CoV-2 cho kết quả dương tính với $76{,}2\%$ các ca thực sự nhiễm virus và kết quả âm tính với $99{,}1\%$ các ca thực sự không nhiễm virus (nguồn https://tapchiyhocvietnam.vn/index.php/vmj/article/view/2124/1921). Giả sử tỉ lệ người thực sự nhiễm virus SARS-CoV-2 trong một cộng  đồng X là~$1\%$. Một người trong cộng đồng đó làm xét nghiệm.
	\choiceTF 
		{\True Tỉ lệ người thực sự  không nhiễm virus SARS-CoV-2 trong cộng đồng X là~$99\%$}
		{Biết rằng người đó thực sự không nhiễm virus, xác suất người  đó có kết quả dương tính bằng $0{,}762$}
		{\True Xác suất người làm xét nghiệm đó có  kết quả dương tính (kết quả làm tròn đến hàng phần nghìn) bằng $0{,}017$}
		{\True Biết người làm xét nghiệm đó nhận được kết quả dương tính, xác suất người đó thực sự nhiễm  virus (kết quả làm tròn đến hàng phần nghìn) bằng $0{,}448$}
	 \loigiai{
		\begin{itemchoice}
			\itemch Tỉ lệ người  không nhiễm virus là~$100\% - 1\% = 99\%$.
			\itemch Sai. Xác suất dương tính  khi không nhiễm virus là~$100\% - 99{,}1\% = 0{,}9\%=0{,}009$.
			\itemch Gọi  $A$ là~biến cố \lq\lq người nhiễm bệnh\rq\rq, $B$ là~biến cố \lq\lq người xét nghiệm có kết quả dương tính\rq\rq.\\
			 $P(A) = 1\% = 0{,}01$ và $P(\overline{A}) = 99\% = 0{,}99$. \\
			$P(B|A) = 76{,}2\% = 0{,}762$ và $P(\overline{B}|\overline{A}) = 1 - 99{,}1\% =  0{,}9\% = 0{,}009$. \\
			Xác suất người có kết quả dương tính là~$$P(B) = P(A)  \cdot P(B|A) + P(\overline{A}) \cdot P(\overline{A}|B) = 0{,}01 \cdot (0{,}762) + 0{,}99 \cdot (0{,}009) \approx 0{,}01653 \approx 0{,}017.$$
			\itemch Xác suất người thực sự nhiễm virus khi kết quả  dương tính là
			$$P(A|B) = \dfrac{P(A) \cdot P(B|A)}{P(B)} = \dfrac{0{,}01 \cdot 0{,}762}{0{,}017} \approx 0{,}448.$$ 
		\end{itemchoice}
	} 
\end{ex}

% nguồn: Lop_12_Dai_so_Chuong_6
\begin{ex}%[ma:21d15458]
	%[2D6H1-2] 
	Một lớp học có $17$ học sinh nam và $24$ học sinh nữ. Cô giáo gọi ngẫu nhiên lần lượt $2$  học sinh (có thứ tự) lên trả lời câu hỏi. Xét các biến cố:\\
	 $A\colon$\lq\lq Lần thứ nhất cô giáo gọi $1$ học sinh nam\rq\rq;\\
	$B\colon$\lq\lq  Lần thứ hai cô giáo gọi $1$ học sinh nữ\rq\rq.
	 \choiceTF
		{$\mathrm{P}(B \mid  A)=0{,}575$}
		{$\mathrm{P}\big(B \mid  \overline{A}\big)=0{,}6$}
		{$\mathrm{P}\big(\overline{B}\mid A\big)=0{,}425$} 
		 {$\mathrm{P}\big(\overline{B}\mid \overline{A}\big)=0{,}4$}
	\loigiai{ 
		\begin{itemchoice}
			\itemch {\bf Sai}.\\ 
			Biến cố $A\cap B\colon$\lq\lq Lần thứ nhất cô giáo gọi  $1$ học sinh nam và lần thứ hai cô giáo gọi $1$ học sinh nữ\rq\rq.\\
			Xác suất xảy ra biến cố $A\cap B$ là $\mathrm{P}\big(A\cap B\big)=\dfrac{17\cdot  24}{\mathrm{C}_{41}^2}$.\\
			Xác suất xảy ra biến cố $A$  là
			$\mathrm{P}\big(A\big)=\dfrac{17\cdot 16+ 17\cdot 24}{\mathrm{C}_{41}^2}$ (tương ứng với hai trường hợp nam  - nam và nam - nữ).\\
			Vậy $\mathrm{P}(B \mid A)=\dfrac{\mathrm{P}\big(A\cap B\big)}{\mathrm{P}(A)}=\dfrac{17\cdot  24}{17 \cdot 16+17 \cdot 24}=\dfrac{3}{5}=0{,}6$.
			\itemch {\bf  Sai}.\\
			Ta có biến  cố $\overline{A}\colon$\lq\lq Lần thứ nhất cô giáo gọi 1 học sinh nữ\rq\rq.\\
			Vậy biến cố $\overline{A}\cap B\colon$\lq\lq Hai học sinh gọi  lên đều là nữ\rq\rq.\\
			 Ta có $\mathrm{P}\big(\overline{A}\cap B\big)=\dfrac{24\cdot 23}{\mathrm{C}_{41}^2}$.\\
			Xác suất xảy ra biến cố $\overline{A}$ là $\mathrm{P}\big( \overline{A}\big)=\dfrac{24\cdot 17+24\cdot 23}{\mathrm{C}_{41}^2}$ (tương ứng với hai trường hợp nữ  - nữ và nữ - nam).\\
			Vậy $\mathrm{P}\big(B \mid \overline{A}\big)=\dfrac{\mathrm{P}\big(\overline{A} \cap B\big)}{\mathrm{P}\big(\overline{A}\big)}=\dfrac{24 \cdot 23}{24 \cdot 17+24 \cdot 23}=\dfrac{23}{40}=0{,}575$. 
			\itemch {\bf Sai}.\\ 
			Ta có biến cố $\overline{B}\colon$\lq\lq Lần thứ hai cô  giáo gọi $1$ học sinh nam\rq\rq.\\
			Vậy biến cố $\overline{B}\cap A\colon$\lq\lq Hai học sinh gọi lên  đều là nam\rq\rq.\\
			Ta  có $\mathrm{P}\big(\overline{B}\cap A\big)=\dfrac{17\cdot 16}{\mathrm{C}_{41}^2}$.\\
			Vậy $\mathrm{P}\big(\overline{B}\mid A\big)=\dfrac{\mathrm{P}(A \cap \overline{B})}{\mathrm{P}(A)}=\dfrac{17 \cdot 16}{17 \cdot  16+17 \cdot 24}=\dfrac{2}{5}=0{,}4$.
			\itemch {\bf Sai}.\\ 
			Vậy biến cố $\overline{A}\cap \overline{B}\colon$\lq\lq Lần thứ nhất cô giáo gọi  $1$ học sinh nữ và lần thứ hai cô giáo gọi $1$ học sinh nam\rq\rq.\\
			Ta có $\mathrm{P}\big(\overline{A}\cap \overline{A}\big)=\dfrac{24\cdot 17}{\mathrm{C}_{41}^2}$.\\ 
			 Vậy $\mathrm{P}\big(\overline{B}\mid \overline{A}\big)=\dfrac{\mathrm{P}\big(\overline{A} \cap \overline{B}\big)}{\mathrm{P}\big(\overline{A}\big)}=\dfrac{24 \cdot 17}{24 \cdot 17+24 \cdot 23}=\dfrac{17}{40}=0{,}425$.
		\end{itemchoice}
	} 
\end{ex}

% nguồn: Lop_12_Dai_so_Chuong_6
\begin{ex}%[ma:0dd47779]
	 %[2D6H1-2]
	Trên một bảng quảng cáo, người ta mắc hai hệ thống bóng đèn. Hệ thống I gồm 2 bóng mắc nối tiếp, hệ thống II gồm 2 bóng mắc song song. Khả năng bị hỏng của mỗi bóng đèn sau 6 giờ thắp sáng liên tục là~$0{,}15$. Biết  tình trạng của mỗi bóng đèn là~độc lập. Xét tính đúng sai của các khẳng định sau.
	\choiceTF 
		{Xác  suất cả hai hệ thống bị hỏng (không sáng) là~$0{,}026$}
		{Xác suất  để hệ thống I bị hỏng (không sáng) là~$0{,}5775$}
		{\True Xác suất để  hệ thống II hoạt động bình thường là~$0{,}9775$}
		{\True Xác  suất để hệ thống II bị hỏng (không sáng) là~$0{,}0225$}
	\loigiai{ 
		Xác suất một bóng đèn  bị hỏng là~$0{,}15$.\\
		Xác  suất một bóng đèn hoạt động bình thường là~$ 1-0{,}15=0{,}85$.
		\begin{itemize}
			\item Sai: Hệ thống I mắc nối tiếp nên hệ thống I hoạt động bình thường khi cả hai  bóng đèn không bị hỏng.Xác suất để hệ thống I hoạt động bình thường là~$0{,}85\cdot 0{,}85$.\\
			Gọi biến cố $A:$ "Hệ thống  I bị hỏng".\\
			$\mathrm{P}(A)=1-0{,}85\cdot 0{,}85=0{,}2775$. 
			\item  Đúng: Hệ thống II mắc song song nên hệ thống II hỏng khi cả hai bóng đèn hỏng.\\
			Gọi biến cố $B:$  "Hệ thống II bị hỏng".\\
			$\mathrm{P}(B)=0{,}15\cdot 0{,}15=0{,}0225$. 
			\item Đúng: Xác suất để  hệ thống II hoạt động bình thường:\\
			$\mathrm{P}\left(\overline{B}\right)=1-0{,}0225=0{,}9775$. 
			\item  Sai: Xác suất cả hai hệ thống bị hỏng:\\
			$\mathrm{P}\left(AB\right)=\mathrm{P}(A)\cdot \mathrm{P}(B)=0{,}2775\cdot  0{,}0225\approx 0{,}062$.
		\end{itemize}
	} 
\end{ex}

% nguồn: Lop_12_Dai_so_Chuong_6
\begin{ex}%[ma:38c1851f]
	 %[2D6H2-4]
	Một trường năng khiếu có 1000 học  sinh. Trong đó có 200 học sinh tham gia câu lạc bộ âm nhạc, và có $ 85\% $ học sinh biết chơi đàn guitar. Ngoài ra, có $ 10\% $ số học sinh không tham gia câu lạc bộ âm nhạc cũng biết chơi đàn guitar. Chọn ngẫu nhiên 1 học sinh của trường.
	\choiceTF 
		{Xác suất chọn được học sinh không  tham gia câu lạc bộ âm nhạc là~$ 0{,}9 $}
		{\True Xác suất chọn được học sinh vừa tham gia câu lạc bộ âm nhạc vừa biết chơi đàn guitar là~$ 0{,}17 $} 
		{\True Xác suất chọn được học sinh biết chơi đàn guitar là~$ 0{,}25 $} 
		{Giả sử học sinh đó biết chơi đàn guitar. Xác  suất chọn được học sinh thuộc câu lạc bộ âm nhạc là~$ 0{,}7 $}
	 \loigiai{
		\begin{itemchoice}
			\itemch $ \mathrm{P}(A)=\dfrac{200}{1000}=0{,}2 \Rightarrow  \mathrm{P}(\overline{A})=0,8 $
			\itemch $ \mathrm{P}(B|A)=0{,}85;\ \mathrm{P}(B|\overline{A})=0,1 $ xác suất cần tìm là~$  \mathrm{P}(AB) $\\
			Theo công thức nhân xác suất, ta  có:\\
			$ \mathrm{P}(AB)=\mathrm{P}(BA)=\mathrm{P}(A).\mathrm{P}(B|A)=0{,}2.0,85=0,17  $
			\itemch Theo công thức  xác suất toàn phần, ta có:\\
			 $ \mathrm{P}(B)=\mathrm{P}(A).\mathrm{P}(B|A)+\mathrm{P}(\overline{A}).\mathrm{P}(B|\overline{A})=0{,}2\cdot0,85+0,8\cdot0,1=0,25 $
			\itemch Theo công thức Bayes, xác suất chọn được học sinh thuộc câu lạc  bộ âm nhạc, biết học sinh đó chơi được đàn guitar là:\\
			$ \mathrm{P}(A|B)=\dfrac{\mathrm{P}(A)\cdot  \mathrm{P}(B|A)}{\mathrm{P}(B)}=\dfrac{0{,}2\cdot0,85}{0{,}25}=0{,}68 $
		\end{itemchoice}
	} 
\end{ex}

% nguồn: Lop_12_Dai_so_Chuong_6
\begin{ex}%[ma:352ea304]
	%[2D6H1-4] 
	Một công ty truyền thông đấu thầu $2$ dự án. Khả năng thắng thầu của dự án 1 là~$0{,}5$ và dự án 2 là~$0{,}6$. Khả năng thắng thầu của 2  dự án là~$0{,}4$. Gọi $A\colon $\lq\lq công ty thắng thầu dự án 1\rq\rq, $B\colon $\lq\lq công ty thắng thầu dự án 2\rq\rq.
	 \choiceTF
		{\True  Xác suất $\mathrm{P}(\overline{A})=0{,}5$ và $\mathrm{P}(\overline{B})=0{,}4$}
		{\True Xác  suất công ty thắng thầu đúng 1 dự án là~$0{,}3$}
		{Biết công ty thắng thầu  dự án 1, xác suất công ty thắng thầu dự án 2 là~$0{,}4$}
		{Biết công ty không thắng thầu dự án 1,  xác suất công ty thắng thầu dự án 2 là~$0{,}8$}
	 \loigiai{
		\begin{itemchoice}
			\itemch  $ \mathrm{P}(A)=0{,}5\Rightarrow \mathrm{P}(\overline{A})=0{,}5$; $\mathrm{P}(B)=0{,}6 \Rightarrow \mathrm{P}(\overline{B})=0{,}4$.
			\itemch Gọi $C\colon$\lq\lq công ty thắng thầu đúng 1 dự  án\rq\rq.\\
			$  \mathrm{P}(C)=\mathrm{P}(A\cap\overline{B})+\mathrm{P}(\overline{A}\cap B)=\mathrm{P}(A)-\mathrm{P}(A\cap B)+\mathrm{P}(B)-\mathrm{P}(A\cap B)$\\
			$=\mathrm{P}(A)+\mathrm{P}(B)-2\mathrm{P}(A\cap B)=0{,}5+0{,}6-2\cdot  0{,}4=0{,}3$.
			\itemch  Gọi $D\colon $\lq\lq công ty thắng thầu dự án 1 biết công ty đã thắng thầu dự án 1\rq\rq.\\
			$  \mathrm{P}(D)=\mathrm{P}(B\mid A)=\dfrac{\mathrm{P}(B\cap A)}{\mathrm{P}(A)}=\dfrac{0{,}4}{0{,}5}=0{,}8$.
			\itemch Gọi $E\colon $\lq\lq công ty thắng thầu dự  án 2 biết công ty không thắng thầu dự án 1\rq\rq.\\
			$\mathrm{P}(E)=\mathrm{P}(B\mid \overline{A})=\dfrac{\mathrm{P}(B\cap\overline{A})}{\mathrm{P}(\overline{A})}=\dfrac{\mathrm{P}(B)-\mathrm{P}(A\cap  B)}{\mathrm{P}(\overline{A})}=\dfrac{0{,}6-0{,}4}{0{,}5}=0{,}4$.
		\end{itemchoice}
	} 
\end{ex}

% nguồn: Lop_12_Dai_so_Chuong_6
\begin{ex}%[ma:898fb24c]
	 %[2D6V1-2]
	Một công ty đấu thầu $2$ dự án. Khả năng thắng thầu của dự án $1$ là~$0{,}4$ và khả năng thắng thầu của dự án $2$ là~$0{,}5$. Khả năng thắng thầu cả  $2$ dự án là~$0{,}3$.\\
	Gọi $A$ là~biến cố: \lq\lq  Thắng thầu dự án $1$\rq\rq.\\
	Gọi $B$ là~biến cố: \lq\lq Thắng thầu dự án $2$ \rq\rq. 
	\choiceTF 
		{$A$ và $B$ là~hai biến cố độc lập} 
		{Xác suất để công ty thắng thầu đúng $1$ dự án bằng  $0{,}7$}
		{\True Xác suất để công ty thắng thầu dự án  $2$ biết công ty thắng thầu dự án $1$ là~$0{,}75$}
		{Xác suất để công ty thắng thầu dự án $2$ biết công ty không thắng thầu dự án  $1$ là~$0{,}25$}
	 \loigiai{
		\begin{itemchoice}
			\itemch Theo giả thiết ta có $\mathrm{P}(A)=0{,}4$; $\mathrm{P}(B)=0{,}5$ và $\mathrm{P}(AB)=0{,}3$ do đó $\mathrm{P}(A)\cdot \mathrm{P}(B)=0{,}4  \cdot 0{,}5=0{,}2 \neq 0{,}3$.\\
			Vậy $A$ và  $B$ là~hai biến cố không độc lập.
			 \itemch Gọi $C$ là~biến cố \lq\lq Thắng thầu đúng $1$ dự án\rq\rq.\\
			Ta có $C=\overline{A} B \cup A \overline{B}$ mà $\overline{A} B$ và $A  \overline{B}$ là~các biến cố xung khắc nên $\mathrm{P}(C)=\mathrm{P}(\overline{A} B)+\mathrm{P}(A \overline{B})$.\\
			Ta có $\mathrm{P}(\overline{A}B)=\mathrm{P}(B)-\mathrm{P}(AB)=0{,}5-0{,}3=0{,}2$ nên $\mathrm{P}(A\overline{B})=\mathrm{P}(A)-\mathrm{P}(A B)=0{,}4-0{,}3=0{,}1$.\\ 
			Vậy $\mathrm{P}(C)=0{,}2+0{,}1=0{,}3$. 
			\itemch Gọi $D$ là~biến cố \lq\lq Thắng thầu dự án $2$  biết công ty thắng thầu dự án $1$\rq\rq.\\
			 Ta có $D=B\mid A$ nên $\mathrm{P}(D)=\mathrm{P}(B \mid A)=\dfrac{\mathrm{P}(AB)}{\mathrm{P}(A)}=\dfrac{0{,}3}{0{,}4}=0{,}75$.
			\itemch Gọi $E$ là~biến cố \lq\lq Thắng thầu dự án $2$ biết công ty không thắng thầu dự  án $1$\rq\rq.\\
			Ta có $E=B \mid  \overline{A}$ do đó $\mathrm{P}(E)=\mathrm{P}(B \mid \overline{A})=\dfrac{\mathrm{P}(\overline{A}B)}{\mathrm{P}(\overline{A})}=\dfrac{\mathrm{P}(B)-\mathrm{P}(A B)}{1-\mathrm{P}(A)}=\dfrac{0{,}5-0{,}3}{1-0{,}4}=\dfrac{0{,}2}{0{,}6}=\dfrac{1}{3}$.
		\end{itemchoice}
	} 
\end{ex}

% nguồn: Lop_12_Dai_so_Chuong_6
\begin{ex}%[ma:7a41712b]
	%[2D6V2-4] 
	Một kho hàng có $85\%$ sản phẩm loại $I$ và $15\%$ sản phẩm loại $II$, trong đó có $1\%$ sản phẩm loại $I$ bị hỏng, $4\%$ sản phẩm loại $II$ bị hỏng. Các sản phẩm có kích thước và hình dạng như nhau. Một khách  hàng chọn ngẫu nhiên $1$ sản phẩm.
	 \choiceTF
		{Xác suất  để không chọn được sản phẩm loại $I$ là~$0{,}85$}
		{\True Xác suất để chọn được sản phẩm không bị hỏng trong số các sản phẩm  loại $I$ là~$0{,}99$}
		{\True Xác suất  chọn được sản phẩm không bị hỏng là~$0{,}9855$}
		{Xác suất  chọn được sản phẩm loại $I$ mà không bị hỏng là~$0{,}95$}
	\loigiai{\begin{itemchoice}
	 \itemch \textbf{Sai}.\\
	Gọi $A$ là~sự kiện  \lq\lq Khách hàng chọn được sản phẩm loại $I$\rq\rq.\\
	Suy ra $\overline{A}$ là~sự kiện \lq\lq Khách hàng  không chọn được sản phẩm loại $I$\rq\rq.\\
	Hay $\overline{A}$ là~sự kiện \lq\lq Khách hàng  chọn được sản phẩm loại $II$\rq\rq.\\
	Theo đề bài ra ta có $P(A) = 0{,}85$ và $P\left( \overline{A} \right)  = 0{,}15$.
	\itemch \textbf{Đúng}.\\ 
	Gọi $B$ là~sự kiện \lq\lq Khách hàng chọn được sản phẩm bị  hỏng\rq\rq.\\
	Ta có, xác suất để chọn được sản phẩm bị hỏng trong  số các sản phẩm loại $I$ là~$P\left(B \mid A \right) = 0{,}01$.\\
	Xác suất để chọn được sản phẩm bị không hỏng trong số các sản phẩm loại $I$ là~$P\left(\overline{B} \mid A \right)  = 0{,}99$.
	\itemch \textbf{Đúng}.\\ 
	Xác  suất chọn được sản phẩm bị hỏng là
	\begin{align*}
		P\left(B \right) &= P(A) \cdot P\left(B \mid A \right)  + P\left(\overline{A} \right) \cdot P\left(B \mid \overline{A} \right) \\
		&= 0{,}85 \cdot 0{,}01 +  0{,}15 \cdot 0{,}04 = 0{,}0145.
	\end{align*}
	Xác suất để chọn được sản phẩm không bị hỏng là~$P\left(\overline{B} \right) = 1  - P(B) = 0{,}9855$.
	\itemch \textbf{Sai}.\\ 
	Xác suất chọn được sản phẩm loại $I$ mà không bị hỏng  là
	\begin{align*}
		P\left(A \mid \overline{B} \right) = \dfrac{P(A) \cdot P\left(\overline{B} \mid A  \right)}{P\left(\overline{B} \right)} = \dfrac{0{,}85 \cdot 0{,}99}{0{,}9855} \approx 0{,}85.
	\end{align*}
\end{itemchoice}
 }
\end{ex}

% nguồn: Lop_12_Dai_so_Chuong_6
\begin{ex}%[ma:d970cb3c]
	 %[2D6H2-2]
	Một cửa hàng có hai loại bóng đèn LED, trong đó có $65 \%$ bóng đèn LED là màu trắng và $35 \%$ bóng đèn LED là màu xanh, các bóng đèn có kích thước như nhau. Các bóng đèn LED màu trắng có tỉ lệ hỏng là $2 \%$ và các bóng đèn LED màu xanh có tỉ lệ hỏng là $3 \%$.  Một khách hàng chọn mua ngẫu nhiên $1$ bóng đèn LED từ cửa hàng. Xét các biến cố:\\
	$A\colon$\lq\lq Khách hàng chọn  được bóng đèn LED màu trắng\rq\rq;\\
	$B\colon$\lq\lq Khách hàng chọn được bóng đèn  LED không hỏng\rq\rq.
	\choiceTF 
		{$\mathrm{P}\big(\overline{A}\big)=0{,}65$} 
		{$\mathrm{P}(B \mid A)=0{,}02$} 
		{$\mathrm{P}\big(B \mid \overline{A}\big)=0{,}3$} 
		{\True  $\mathrm{P}(B)=0{,}9\,765$}
	\loigiai{ 
		\begin{itemchoice}
			\itemch {\bf Sai}.\\ 
			Xác suất khách hàng chọn được bóng đèn  LED màu trắng là $\mathrm{P}(A)=0{,}65$.\\
			Biến cố  $\overline{A}\colon$\lq\lq Khách hàng chọn được bóng đèn LED màu trắng\rq\rq.\\
			Xác  suất xảy ra biến cố $\overline{A}$ là
			 $\mathrm{P}\big(A\big)=1-P(A)=1-0{,}65=0{,}35$.
			\itemch {\bf Sai}.\\ 
			Ta  có $\mathrm{P}(B \mid A)=1-\mathrm{P}\big(\overline{B}\mid A\big)=1-0{,}02=0{,}98$.
			\itemch {\bf Sai}.\\ Ta có $\mathrm{P}\big(B \mid  \overline{A}\big)=1-\mathrm{P}\big(\overline{B} \mid \overline{A}\big)=1-0{,}03=0{,}97$.
			\itemch {\bf Đúng}.\\ Theo công thức  xác suất toàn phần ta có \[\mathrm{P}(B)=\mathrm{P}(A)\cdot \mathrm{P}(B \mid A)+\mathrm{P}\big(\overline{A}\big)\cdot \mathrm{P}\big(B\mid \overline{A}\big)=0{,}65\cdot 0{,}98+0{,}35\cdot 0{,}97=0{,}9765.\]
		\end{itemchoice}
	 }
\end{ex}

% nguồn: Lop_12_Dai_so_Chuong_6
\begin{ex}%[ma:87bde602]
	 %[2D6V2-3]
	Khi tìm hiểu về các hồ sơ bệnh án liên quan đến  các bệnh nhân bị bỏng, các bác sĩ đã rút ra $50$ hồ sơ bệnh án của $50$ bệnh nhân. Trong đó có $35$ bệnh nhân bị bỏng do nhiệt và $15$ bệnh nhân bị bỏng do hóa chất. Trong quá trình diễn biến của bệnh thì trong số $35$ bệnh nhân bị bỏng nhiệt có $30\%$ bệnh nhân bị biến chứng và trong số $15$ bệnh nhân bị bỏng do hóa chất thì có $50\%$ bệnh nhân bị biến chứng.\\
	Gọi $A$ là~biến cố: \lq \lq Hồ sơ rút  ra bệnh nhân bị bỏng do nhiệt\rq\rq;\\
	$B$ là~biến cố: \lq\lq Hồ sơ rút ra  bệnh nhân bị bỏng do hóa chất \rq\rq;\\
	 $C$ là~biến cố: \lq\lq Hồ sơ rút ra bệnh nhân bị biến chứng \rq\rq.
	\choiceTF 
		 {\True $\mathrm{P}(B) = \dfrac{3}{10}$, $\mathrm{P}(\overline{B}) = \dfrac{7}{10}$}
		{\True Xác suất có  điều kiện $\mathrm{P}(C\mid A)=0{,}3 $}
		{$\mathrm{P}(C) = 0{,}85$} 
		{Trong số các bệnh nhân bị biến chứng có $25\%$ bệnh nhân bị bỏng do nhiệt (kết quả tính theo phần trăm được làm tròn đến hàng đơn  vị)}
	\loigiai{ 
	\begin{itemchoice}
		\itemch $\mathrm{P}(B) =\dfrac{n(B)}{n\left(\Omega\right)}= \dfrac{15}{50} = \dfrac{3}{10}$,  $\mathrm{P}(\overline{B}) =1-\mathrm{P}(B)= \dfrac{7}{10}$.
		 \itemch
		Khi biến cố $A$ xảy ra tức là~bệnh nhân bị bỏng do nhiệt thì xác suất bệnh nhân  bị biến chứng là~$0{,}3$.
		\itemch 
		Ta có:  $\mathrm{P}(A)=0{,}7,$ $\mathrm{P}(\overline{A})=0{,}3 $, $\mathrm{P}\left(C\mid A\right)=0{,}3$; $\mathrm{P}\left(C\mid \overline{A}\right)=0{,}5$.\\
		Áp dụng công thức xác suất  toàn phần ta có\\
		$\mathrm{P}(C)=\mathrm{P}(A)\cdot \mathrm{P}\left(C\mid A\right)+\mathrm{P}\left(\overline{A}\right)\cdot \mathrm{P}\left(C\mid \overline{A}\right)=0{,}7\cdot 0{,}3+0{,}3\cdot  0{,}5=0{,}36$.
		\itemch 
		Nếu hồ sơ rút ra bệnh nhân bị biến chứng thì xác suất để bệnh nhân bị bỏng do nhiệt theo công thức  Bayes ta có\\
		$\mathrm{P}\left(A\mid C\right)=\dfrac {\mathrm{P}(A)\cdot \mathrm{P}\left(C\mid A\right)}{\mathrm{P}(C)}=\dfrac {0{,}7\cdot 0{,}3}{0{,}36}\approx  0{,}58$.\\
		Vậy trong số các bệnh nhân bị biến chứng có $58\%$  bệnh nhân bị bỏng do nhiệt.
	\end{itemchoice}
	 }
\end{ex}

% nguồn: Lop_12_Dai_so_Chuong_6
\begin{ex}%[ma:7ff41bae]
	 %[2D6H1-2]
	Trong một hộp có  $18$ quả bóng bàn loại I và $2$ quả bóng bàn loại II, các quả bóng bàn có hình dạng và kích thước như nhau. Một học sinh lấy ngẫu nhiên lần lượt $2$ quả bóng bàn (lấy không hoàn lại) trong hộp.
	\choiceTF 
		{Xác suất để lần thứ nhất lấy được quả bóng bàn  loại II là $\dfrac{9}{10}$}
		{\True  Xác suất để lần thứ hai lấy được quả bóng bàn loại II, biết lần thứ nhất lấy được quả bóng bàn loại II là $\dfrac{1}{19}$}
		{Xác suất  để cả hai lần đều lấy được quả bóng bàn loại II là $\dfrac{9}{190}$}
		{\True Xác suất để ít nhất $1$ lần lấy được quả  bóng bàn loại I là $\dfrac{189}{190}$}
	\loigiai{ 
		\begin{itemchoice}
			\itemch {\bf Sai}.\\ Gọi $A$ là biến cố \lq \lq Lần thứ nhất lấy được quả bóng  bàn loại II\rq \rq.\\
			Khi  đó $\mathrm{P}(A)=\dfrac{2}{20}=\dfrac{1}{10}$.
			\itemch {\bf Đúng}.\\ Gọi $B$ là biến cố \lq \lq Lần thứ hai lấy  được quả bóng bàn loại II\rq \rq.
			Vì lần thứ nhất đã lấy được $1$ quả bóng bàn loại II nên chỉ còn $1$ quả bóng bàn loại II trong  hộp. Do đó xác suất lấy được quả bóng bàn loại II, biết lần thứ nhất đã lấy được quả bóng bàn loại II là $\mathrm{P}(B \mid A)=\dfrac{1}{19}$.
			\itemch {\bf Sai}.\\ Gọi $C$ là biến cố \lq \lq Cả hai lần đều lấy  được quả bóng bàn loại II\rq \rq.
			Khi đó $C = A \cap B$. Do đó xác suất  cần tìm là
			\[\mathrm{P}(C)=\mathrm{P}(A  \cap B)= \mathrm{P}(B \mid A) \cdot \mathrm{P}(A)= \dfrac{1}{19} \cdot \dfrac{1}{10}=\dfrac{1}{190}.\]
			\itemch {\bf  Đúng}.\\ Ta có biến cố đối của biến cố $C$ là $\overline{C}$: \lq \lq Có ít nhất $1$ lần lấy được quả bóng bàn loại I\rq \rq.
			Do đó  $\mathrm{P}\big(\overline{C}\big)=1-\mathrm{P}(C)=1-\dfrac{1}{190}=\dfrac{189}{190}$.
		\end{itemchoice}
	} 
\end{ex}

% nguồn: Lop_12_Dai_so_Chuong_6
\begin{ex}%[ma:e24dac2b]
	 %[2D6H1-4]
	Năm $2001$, Cộng đồng Châu Âu có làm một đợt kiểm tra rất rộng rãi các con bò để phát hiện những con bò bị bệnh bò điên. Người  ta tiến hành một loại xét nghiệm và cho kết quả như sau: Khi con bò bị bệnh bò điên thì xác suất để ra phản ứng dương tính trong xét nghiệm là~$70\%$; còn khi con bò không bị bệnh thì xác suất để xảy ra phản ứng dương tính trong xét nghiệm đó là~$10\%$. Biết rằng tỉ lệ bò bị mắc bệnh bò điên ở Hà La là~$1{,}3$ con trên $100\,000$ con. Gọi $X$ là~biến cố một con bò bị bệnh bò điên, $Y$ là~biến cố một con bò phản ứng dương tính với xét nghiệm.
	\choiceTF 
		{\True $P\left(X \right) = 13 \cdot  10^{-6}$}
		{$P\left(  Y\big| X\right) = 0{,}07$}
		{\True $P\left(Y  \big| \overline{X} \right) = 0{,}1$}
		{Xác suất đê kiểm tra trúng một con bò bị bệnh bò  điên và có phản ứng dương tính với xét nghiệm là~$91 \cdot 10^{-8}$}
	\loigiai{\begin{itemchoice}
	 \itemch \textbf{Đúng}.\\
	Tỉ lệ bò bị mắc bệnh bò  điên ở Hà Lan là~$1{,}3$ con trên $100\, 000$ con nghĩa là~$P\left(X \right) = 13 \cdot 10^{-6}$.
	\itemch \textbf{Sai}.\\ 
	Khi con bò bị bệnh bò  điên, thì xác suất để ra phản ứng dương tính trong xét nghiệm là~$70\%$, nghĩa là~$P\left( Y\big| X\right) = 0{,}7$.
	 \itemch \textbf{Đúng}.\\
	Khi con bò không bị bệnh, thì xác suất để ra phản ứng dương tính trong xét nghiệm là~$10\%$, nghĩa là~$P\left(Y \big|  \overline{X} \right) = 0{,}1$.
	 \itemch \textbf{Sai}.\\
	Ta có  $P\left(Y \cap X \right) = P\left(X \right) \cdot P\left( Y\big| X\right) = 13 \cdot 10^{-6} \cdot 0{,}7 = 91 \cdot 10^{-7}$.
\end{itemchoice}
 }
\end{ex}

% nguồn: Lop_12_Dai_so_Chuong_6
\begin{ex}%[ma:e4143f38]
	 %[2D6V2-3]
	Trong một xưởng  máy, số linh kiện do cơ sở I sản xuất chiếm $61 \%$, số linh kiện do cơ sở II sản xuất chiếm $39 \%$. Tỉ lệ linh kiện đạt tiêu chuẩn của cơ sở I, cơ sở II lần lượt là~$93\%$, $82\%$. Kiểm tra ngẫu nhiên một linh kiện ở xưởng máy. Xét các biến cố\\
	$A_1 \colon$ \lq\lq  Linh kiện được kiểm tra do cơ sở I sản xuất\rq\rq;\\
	$A_2 \colon$ \lq\lq Linh kiện được kiểm tra  do cơ sở II sản xuất\rq\rq;\\
	$B \colon$ \lq\lq  Linh kiện được kiểm tra đạt tiêu chuẩn\rq\rq.
	\choiceTF 
		{\True  Xác suất $\mathrm{P}\left(A_1\right)=0{,}61$}
		{\True Xác suất có điều kiện $\mathrm{P}\left(B  \mid A_2\right)=0{,}82$}
		{\True Xác suất  $\mathrm{P}(B)=0{,}8871$}
		{Xác suất có điều kiện  $\mathrm{P}\left(A_1 \mid B\right)=0{,}55$}
	 \loigiai{
		\begin{itemchoice}
			\itemch Ta có  $\mathrm{P}\left(A_1\right)=0{,}61$.
			\itemch Ta có $\mathrm{P}\left(B  \mid A_2\right)=0{,}82$.
			\itemch Ta có  $\mathrm{P}\left(A_1\right)=0{,}61$; $\mathrm{P}\left(A_2\right)=0{,}39$; $\mathrm{P}\left(B \mid A_1\right)=0{,}93$; $\mathrm{P}\left(B \mid A_2\right)=0{,}82$.\\
			Theo công  thức xác suất toàn phần, ta có\\
			$\mathrm{P}(B)=\mathrm{P}\left(A_1\right) \cdot \mathrm{P}\left(B \mid A_1\right)+\mathrm{P}\left(A_2\right) \cdot  \mathrm{P}\left(B \mid A_2\right)=0{,}61 \cdot 0{,}93+0{,}39 \cdot 0{,}82=0{,}8871$.
			\itemch Theo công  thức Bayes, ta có $\mathrm{P}\left(A_1 \mid B\right)=\dfrac{\mathrm{P}\left(A_1\right) \cdot \mathrm{P}\left(B \mid A_1\right)}{\mathrm{P}(B)}=\dfrac{0{,}61 \cdot 0{,}93}{0{,}8871} \approx 0{,}64$.
		\end{itemchoice}
	 }
\end{ex}

% nguồn: Lop_12_Dai_so_Chuong_6
\begin{ex}%[ma:35a7fbfe]
	 %[2D6H1-2]
	Lớp 12A có $40$ học sinh, trong đó có $25$ học sinh tham gia câu lạc bộ Tiếng Anh, $16$ học sinh tham gia câu lạc bộ Toán, $12$ học sinh vừa tham gia câu lạc bộ Tiếng Anh vừa tham gia câu lạc bộ Toán. Chọn ngẫu nhiên $1$ học sinh. Xét các biến cố  sau:
	$A$:  \lq \lq Học sinh được chọn tham gia câu lạc bộ Tiếng Anh\rq \rq;
	$B$: \lq \lq Học sinh được chọn tham gia câu lạc bộ Toán\rq  \rq.
	\choiceTF 
		{$\mathrm{P}(A)=0{,}4$} 
		{$\mathrm{P}(B)=0{,}625$} 
		{\True  $\mathrm{P}(A \mid B)=0{,}75$}
		{\True$\mathrm{P}(B  \mid A)=0{,}48$}
	 \loigiai{
		\begin{itemchoice}
			\itemch {\bf Sai}.\\ Chọn  ngẫu nhiên $1$ trong $40$ học sinh có $\mathrm{C}_{40}^1$ cách. \\Do đó $n(\Omega)=\mathrm{C}_{40}^1=40$.\\
			Vì có $25$ học sinh tham  gia câu lạc bộ Tiếng Anh nên $n(A)=25$.\\
			Vậy  $\mathrm{P}(A)=\dfrac{n(A)}{n(\Omega)}=\dfrac{25}{40}=0{,}625$.
			 \itemch {\bf Sai}.\\ Vì có $16$ học sinh tham gia câu lạc bộ Toán nên $n(B)~=~16$.\\
			Vậy $\mathrm{P}(B)=\dfrac{n(B)}{n(\Omega)}=\dfrac{16}{40}=0{,}4$. 
			\itemch {\bf Đúng}.\\ Vì có $12$ học sinh vừa tham gia câu lạc bộ Tiếng Anh  vừa tham gia câu lạc bộ Toán nên $n(A \cap B)=12$.\\
			Do đó  $\mathrm{P}(A \mid B)=\dfrac{n(A \cap B)}{n(B)}=\dfrac{12}{16}=0{,}75$.
			\itemch {\bf Đúng}.\\ Số học sinh vừa tham gia câu lạc bộ Tiếng Anh vừa tham gia  câu lạc bộ Toán là $12$, số học sinh tham gia câu lạc bộ Tiếng Anh là $25$ nên $\mathrm{P}(B \mid A)=\dfrac{n(A \cap B)}{n(A)}=\dfrac{12}{25}=0{,}48$.
		\end{itemchoice}
	 }
\end{ex}

% nguồn: Lop_12_Dai_so_Chuong_6
\begin{ex}%[ma:044cf378]
	%[2D6V2-2] 
	Một công ty đấu thầu  $2$ dự án. Khả năng thắng thầu của dự án $1$ là~$0{,}4$ và khả năng thắng thầu của dự án $2$ là~$0{,}5$. Khả năng thắng thầu cả $2$ dự án là~$0{,}3$.\\
	Gọi $A$ là~biến cố: \lq\lq Thắng thầu dự án $1$\rq\rq. 
	Gọi $B$ là~biến cố: \lq\lq Thắng thầu  dự án $2$\rq\rq.
	 \choiceTF
		{$A$  và $B$ là~hai biến cố độc lập}
		{\True Xác suất để  công ty thắng thầu đúng $1$ dự án bằng $0{,}7$}
		{Xác suất để công ty thắng thầu dự án $2$ biết công  ty thắng thầu dự án $1$ là~$0{,}75$}
		{\True Xác suất để công ty thắng thầu dự án $2$ biết công ty không  thắng thầu dự án $1$ là~$0{,}25$}
	\loigiai{\begin{itemchoice}
	 \itemch \textbf{Sai}.\\
	Theo giả thiết suy ra: $\mathrm{P}(A) = 0{,}4$, $\mathrm{P}(B) =  0{,}5$ và $\mathrm{P}(AB) = 0{,}3$.\\
	Có $\mathrm{P}(A) \cdot \mathrm{P}(B) = 0{,}4 \cdot  0{,}5 = 0{,}2 \neq 0{,}3 \Rightarrow A$ và $B$ là~hai biến cố không độc lập.
	\itemch  \textbf{Đúng}.\\
	Gọi $C$ là~biến  cố: \lq\lq Thắng thầu đúng $1$ dự án\rq\rq \,$\Rightarrow C = \overline{A}B \cup A\overline{B}$ mà $\overline{A}B$ và $A\overline{B}$ là~các biến cố xung khắc $\Rightarrow P(C) = \mathrm{P}(\overline{A}B) + \mathrm{P}(A\overline{B})$.
	 \[
	\mathrm{P}(\overline{A}B) = \mathrm{P}(B)-\mathrm{P}(AB) = 0{,}5 - 0{,}3 = 0{,}2; 
	\] 
	\[ 
	\mathrm{P}(A\overline{B}) = \mathrm{P}(A) - \mathrm{P}(AB) = 0{,}4 - 0{,}3 =  0{,}1.
	\] 
	Vậy  $
	\mathrm{P}(C) =  0{,}2 + 0{,}1 = 0{,}3.$
	 \itemch \textbf{Đúng}.\\
	Gọi $D$ là~biến  cố: \lq\lq Thắng thầu dự án $2$ biết công ty thắng thầu dự án $1$\rq\rq\\
	$\Rightarrow D  = B|A$.
	Khi  đó:
	 \[
	\mathrm{P}(D) = \mathrm{P}(B|A) = \frac{\mathrm{P}(AB)}{\mathrm{P}(A)} =  \dfrac{0{,}3}{0{,}4} = 0{,}75.
	\] 
	 \itemch \textbf{Sai}.\\
	Gọi $E$ là~biến cố: \lq\lq Thắng thầu dự án $2$ biết công  ty không thắng thầu dự án $1$ \rq\rq $\Rightarrow E = B|\overline{A}$.\\
	 Khi đó:
	 \[
	\mathrm{P}(E) = \mathrm{P}(B|\overline{A}) = \dfrac{\mathrm{P}(\overline{A}B)}{\mathrm{P}(\overline{A})}  = \frac{\mathrm{P}(B) - P(AB)}{1 - \mathrm{P}(A)} = \dfrac{0{,}5 - 0{,}3}{1 - 0{,}4} = \frac{0{,}2}{0{,}6} = \dfrac{1}{3}.
	 \]
\end{itemchoice}
 }
\end{ex}


%%%%%%%%%%%%%%%% TRẢ LỜI NGẮN %%%%%%%%%%%%%%%%

% nguồn: Lop_12_Dai_so_Chuong_6
\begin{bt}%[ma:a6ed14ab]
	 %[2D6H2-3]
	Trong một đợt kiểm tra sức khỏe, có một loại bệnh X mà tỉ lệ người mắc bệnh là $0{,}2\%$ và một loại xét nghiệm Y mà ai mắc bệnh X khi xét nghiệm Y cũng có phản ứng dương tính. Tuy nhiên, có $6\%$ những người không bị bệnh X lại có phản ứng dương tính với xét nghiệm Y. Chọn ngẫu nhiêu $1$ người trong đợt kiểm tra  sức khỏe đó. Giả sử người đó có phản ứng dương tính với xét nghiệm Y. Xác suất người đó bị mắc bệnh X là bao nhiêu (làm tròn kết quả đến hàng phần~trăm)?
	 \shortans{0,03}
	 \loigiai{
		Gọi $A$ là  biến cố \lq \lq Người được chọn bị mắc bệnh X\rq \rq.\\
		Khi  đó $\mathrm{P}(A)=0{,}02$ và $\mathrm{P}\big(\overline{A}\big)=1-0{,}02=0{,}998$.\\
		Gọi $B$ là  biến cố \lq \lq Người được chọn có phản ứng dương tính với xét nghiệm Y\rq \rq.\\
		Theo đề bài, ta có $\mathrm{P}(B \mid A)=1$, $\mathrm{P} \big(B\mid\overline{A}\big)=0{,}06$.\\ 
		Áp dụng  công thức Bayes, ta có
		\[\mathrm{P}\left( {A \mid B} \right) = \dfrac{{\mathrm{P}\left( A \right) \cdot \mathrm{P}\left( {B \mid A} \right)}}{{\mathrm{P}\left( A \right) \cdot \mathrm{P}\left( {B \mid A} \right) + \mathrm{P}\left( {\overline A } \right) \cdot \mathrm{P}\left( {B \mid  \overline A } \right)}} = \dfrac{{0{,}002 \cdot 1}}{{0{,}002 \cdot 1 + 0{,}998 \cdot 0{,}06}}=0{,}03.\]
		Vậy  nếu người đó có phản ứng dương tính với xét nghiệm Y thì xác suất người đó mắc bệnh X xấp xỉ khoảng $0{,}03$.
	 }
\end{bt}

% nguồn: Lop_12_Dai_so_Chuong_6
\begin{bt}%[ma:55964618]
	 %[2D6H2-2]
	Một cuộc thi khoa học có $36$ bộ câu hỏi, trong đó  có $20$ bộ câu hỏi về chủ đề tự nhiên và $16$ bộ câu hỏi về chủ đề xã hội. Bạn An lấy ngẫu nhiên $1$ bộ câu hỏi (lấy không hoàn lại), sau đó bạn Bình lấy ngẫu nhiên $1$ bộ câu hỏi. Xác suất bạn Bình lấy được bộ câu hỏi về chủ đề xã hội bằng $\dfrac{a}{b}$ với $\dfrac{a}{b}$ là phân số tối giản. Giá trị $a+b$ bằng bao nhiêu?
	\shortans{13} 
	 \loigiai{
		 Gọi $A$ là biến cố \lq \lq Bạn An lấy được bộ câu hỏi về chủ đề xã hội\rq \rq.\\
		Khi đó  $\mathrm{P}(A)=\dfrac{16}{36}=\dfrac{4}{9}$ và $\mathrm{P}\big(\overline{A}\big)=1-\mathrm{P}(A)=1-\dfrac{4}{9}=\dfrac{5}{9}$.\\
		Gọi $B$ là biến cố \lq \lq Bạn Bình lấy được bộ  câu hỏi về chủ đề xã hội\rq \rq.\\
		Nếu bạn An chọn được một bộ câu hỏi về chủ đề tự nhiên thì sau đó còn  $35$ bộ câu hỏi, trong đó có $16$ bộ câu hỏi về chủ đề xã hội, suy ra $\mathrm{P} (B\mid A)=\dfrac{16}{35}$.\\
		Nếu bạn An chọn được một bộ câu hỏi về  chủ đề xã hội thì sau đó còn $35$ bộ câu hỏi, trong đó có $15$ bộ câu hỏi về chủ để xã hội, suy ra $\mathrm{P} \big(B\mid\overline{A}\big)=\dfrac{15}{35}$.\\
		Theo công thức xác suất toàn phần, ta  có
		\[\mathrm{P}(B)=\mathrm{P}(B \mid A) \cdot \mathrm{P}(A)+ \mathrm{P}  \big(B\mid\overline{A}\big) \cdot \mathrm{P}\big(\overline{A}\big)=\dfrac{15}{35} \cdot \dfrac{4}{9}+ \dfrac{16}{35} \cdot \dfrac{5}{9}=\dfrac{4}{9}.\]
		Suy ra  $a=4$; $b=9$.\\
		Vậy  $a+b=13$.
	} 
\end{bt}

% nguồn: Lop_12_Dai_so_Chuong_6
\begin{bt}%[ma:a08a5711]
	%[2D6V1-1] 
	Trong kỳ thi tốt nghiệp THPT năm $2025$ có môn thi bắt buộc là môn Toán. Cấu trúc đề thi môn toán gồm $22$ câu với ba phần như sau: 
	\begin{itemize}
		\item Phần $1$: $12$ câu  hỏi ở dạng thức trắc nghiệm nhiều lựa chọn, mỗi câu cho $04$ phương án trả lời trong đó có $01$ phương án là đáp án đúng;
		\item Phần $2$: $04$ câu hỏi ở  dạng thức trắc nghiệm dạng Đúng/Sai. Mỗi câu hỏi có $04$ ý, tại mỗi ý thí sinh lựa chọn đúng hoặc sai;
		\item Phần $3$: $06$ câu hỏi ở dạng thức trắc nghiệm dạng trả lời ngắn. Thí sinh ghi và  tô vào các ô tương ứng với đáp án của mình.
	\end{itemize}
	\textbf{Phần $1$}: Mỗi câu  trả lời đúng được $0{,}25$ điểm.\\
	\textbf{Phần $2$}: 
	\begin{itemize}
		\item  Chỉ lựa chọn chính xác $01$ ý trong $01$ câu hỏi được $0{,}1$ điểm;
		\item Chỉ lựa chọn  chính xác $02$ ý trong $01$ câu hỏi được $0{,}25$ điểm;
		\item Chỉ lựa chọn chính xác $03$ ý trong $01$ câu  hỏi được $0{,}5$ điểm;
		\item Chính xác cả $04$ ý trong $01$ câu hỏi  được $1$ điểm.
	\end{itemize}
	\textbf{Phần $3$}: Mỗi câu trả lời đúng được  $0{,}5$ điểm.\\
	Giả sử một thí sinh do lười học nên không làm được môn toán. Phần $3$ thí sinh không trả lời đúng câu nào. Phần $1$ và phần $2$ thí sinh tô bừa các đáp án. Nếu thí sinh tô phần $2$ được tổng $0{,}55$ điểm. Khi đó xác suất của  $12$ câu tô bừa ở phần $1$ để tổng điểm bài làm môn toán của thí sinh được trên $3$ điểm là $x\cdot 10^{-8}$. Tổng các chữ số phần nguyên của $x$ bằng bao~nhiêu?
	\shortans{17} 
	\loigiai{ 
		Để tổng điểm bài làm môn toán của thí sinh được trên $3$ điểm thì số điểm ở phần $1$ sẽ lớn hơn $2{,}45$ điểm. Vì mỗi câu trả lời đúng phần $1$ được $0{,}25$ điểm nên số câu trả lời đúng sẽ lớn hơn hoặc  bằng $10$.
		\begin{itemize}
			\item Trường hợp $1$: Số  câu đúng ở phần $1$ là $10$ câu. Khi đó xác suất là $P_1=0{,}25^{10}\cdot 0{,}75^2\cdot C_{12}^{10}$.
			\item Trường hợp $2$: Số câu đúng ở phần $1$ là $11$ câu. Khi đó xác suất là $P_2=0{,}25^{11}\cdot 0{,}75\cdot C_{12}^{11}$. 
			\item Trường hợp $3$: Số câu đúng ở phần $1$  là $12$ câu. Khi đó xác suất là $P_3=0{,}25^{12}$.
		\end{itemize}
		Vậy  xác suất cần tìm là $P=P_1+P_2+P_3 \approx 3761{,}053085\cdot10^{-8}$.\\
		Do đó  phần nguyên của $x$ là $3761$.\\
		Vậy tổng các chữ số  phần nguyên của $x$ bằng 17.
	} 
\end{bt}

% nguồn: Lop_12_Dai_so_Chuong_6
\begin{bt}%[ma:4130067a]
	 %[2D6V2-3]
	Trường Bình Phúc có $20  \%$ học sinh tham gia câu lạc bộ âm nhạc, trong số học sinh đó có $85 \%$ học sinh biết chơi đàn guitar. Ngoài ra, có $10 \%$ số học sinh không tham gia câu lạc bộ âm nhạc cũng biết chơi đàn guitar. Chọn ngẫu nhiên $1$ học sinh của trường. Giả sử học sinh đó biết chơi đàn guitar. Xác suất chọn được học sinh thuộc câu lạc bộ âm nhạc là bao nhiêu?
	\shortans{0,68} 
	\loigiai{Xét các biến cố $A$: \lq \lq Chọn được học sinh thuộc câu lạc bộ âm nhạc\rq  \rq; \\
		$B$: \lq \lq Chọn được học  sinh biết chơi đàn guitar\rq \rq.\\
		Khi  đó, $\mathrm{P}(A)=0{,}2$; $\mathrm{P}\big(\overline{A}\big)=0{,}8$; $\mathrm{P}(B \mid A)=0{,}85$; $\mathrm{P}\big(B \mid \overline{A}\big)=0{,}1$.\\
		Theo công thức  xác suất toàn phần, ta có
		\[ 
		\mathrm{P}(B)=\mathrm{P}(A) \cdot \mathrm{P}(B \mid A)+\mathrm{P}\big(\overline{A}\big) \cdot \mathrm{P}\big(B \mid \overline{A}\big)=0{,}2  \cdot 0{,}85+0{,}8 \cdot 0{,}1=0{,}25 .
		\] 
		Theo công thức Bayes, xác suất chọn được học sinh thuộc câu lạc bộ âm nhạc, biết học sinh đó chơi  được đàn guitar, là
		 \[
		\mathrm{P}(A \mid B)=\frac{\mathrm{P}(A) \cdot  \mathrm{P}(B \mid A)}{\mathrm{P}(B)}=\frac{0{,}2 \cdot 0{,}85}{0{,}25}=0{,}68.
		 \]
	} 
\end{bt}

% nguồn: Lop_12_Dai_so_Chuong_6
\begin{bt}%[ma:3ffaed73]
	%[2D6V1-2] 
	Một căn bệnh có $1\%$ dân số mắc phải. Một phương pháp chuẩn đoán được phát triển có tỷ lệ chính xác là $99\%$. Với những người bị  bệnh, phương pháp này sẽ đưa ra kết quả dương tính $99\%$ số trường hợp. Với người không mắc bệnh, phương pháp này cũng chuẩn đoán đúng $99$ trong $100$ trường hợp. Nếu một người kiểm tra và kết quả là dương tính (bị bệnh), xác suất để người đó thực sự bị bệnh là bao~nhiêu?
	\shortans{0,5} 
	\loigiai{ 
		Gọi  $A\colon $\lq\lq người đó mắc bệnh\rq\rq.\\
		Gọi $B\colon $\lq\lq kết quả kiểm  tra người đó là dương tính (bị bệnh)\rq\rq.\\
		Ta cần tính  $\mathrm{P}\left(A\mid B\right)$.\\
		Với $\mathrm{P}\left(A\mid B\right)=\dfrac{\mathrm{P}(A)\cdot \mathrm{P}\left(B\mid  A\right)}{\mathrm{P}(A)\cdot \mathrm{P}\left(B\mid A\right)+\mathrm{P}\left(\overline{A}\right)\cdot \mathrm{P}\left(B\mid \overline{A}\right)}$.\\
		Xác suất để người đó mắc bệnh khi chưa kiểm tra  $\mathrm{P}(A)=1\%=0{,}01$.\\
		 Do đó xác suất để người đó không mắc bệnh khi chưa kiểm tra $\mathrm{P}\left(\overline{A}\right)=1-0{,}01=0{,}99$.\\
		Xác suất kết quả dương tính nếu người đó mắc  bệnh là $\mathrm{P}\left(B\mid A\right)=99\%=0{,}99$.\\
		Xác suất kết quả dương tính nếu người đó không mắc bệnh là  $\mathrm{P}\left(B\mid \overline{A}\right)=1-0{,}99=0{,}01$.\\
		$\mathrm{P}\left(A\mid B\right)=\dfrac{\mathrm{P}(A)\cdot \mathrm{P}\left(B\mid A\right)}{\mathrm{P}(A)\cdot \mathrm{P}\left(B\mid A\right)+\mathrm{P}\left(\overline{A}\right)\cdot P\left(B\mid \overline{A}\right)}=\dfrac{0{,}01\cdot 0{,}99}{0{,}01\cdot 0{,}99+0{,}99\cdot  0{,}01}=0{,}5$.\\
		Xác suất để người đó mắc bệnh nếu kết  quả kiểm tra người đó là dương tính là $0{,}5$.
	 }
\end{bt}

% nguồn: Lop_12_Dai_so_Chuong_6
\begin{bt}%[ma:f9c5442c]
	 %[2D6C2-4]
	Mộthộp chứa $10$ viên bi xanh và $5$ viên bi đỏ. Bạn An lấy ra ngẫu nhiên $1$ viên bi từ hộp,  xem màu, rồi bỏ ra ngoài. Nếu viên bi An lấy ra có màu xanh, bạn Bình sẽ lấy ra ngẫu nhiên $2$ viên bi từ hộp; còn nếu viên bi An lấy ra có màu đỏ, bạn Bình sẽ lấy ra ngẫu nhiên $3$ viên bi từ hộp. Tính xác suất để An lấy được viên bi màu xanh, biết rằng tất cả các viên bi được hai bạn chọn ra đều có đủ cả hai màu.
	\shortans{0,55} 
	\loigiai{ 
		Gọi $A$ là biến cố \lq\lq An lấy được  viên bi xanh\rq\rq.\\
		Khi đó $\overline{A}$ là biến cố \lq\lq An lấy được viên bi màu  đỏ\rq\rq.\\
		Gọi $B$ là biến cố \lq\lq tất  cả các viên bi được hai bạn chọn ra đều có đủ hai màu\rq\rq.\\
		Trường hợp $1$. An lấy  được viên bi màu xanh.\\
		Khi đó hộp còn $9$ viên bi màu xanh và $5$  viên bi màu đỏ.\\
		Để có đủ $2$ màu thì  có các trường hợp là
		\begin{itemize}
			\item Bình lấy hai viên  bi đỏ.
			\item Bình lấy $1$ viên bi xanh  và $1$ viên bi đỏ.
		\end{itemize}
		Suy ra $\mathrm{P}\left(B \mid A \right) = \dfrac{\mathrm{C}_{5}^{2} + \mathrm{C}_{9}^{1} \cdot  \mathrm{C}_{5}^{1}}{\mathrm{C}_{14}^{2}} = \dfrac{55}{91}$.\\
		Trường hợp $2$. An lấy được viên bi đỏ.\\ 
		 Khi đó hộp còn $10$ viên bi xanh và $4$ viên bi đỏ.\\
		Để  có đủ $2$ màu thì có các trường hợp là
		\begin{itemize}
			\item Bình lấy $3$ viên bi  xanh.
			\item Bình lấy $2$ viên bi xanh và $1$ viên  bi đỏ.
			\item Bình lấy $1$  viên bi xanh và $2$ viên bi đỏ.
		\end{itemize}
		Suy ra $\mathrm{P}\left(B \mid \overline{A} \right) = \dfrac{\mathrm{C}_{10}^{3} +  \mathrm{C}_{10}^{2} \cdot \mathrm{C}_{4}^{1} + \mathrm{C}_{10}^{1} \cdot \mathrm{C}_{4}^{2}}{\mathrm{C}_{3}^{14}} = \dfrac{90}{91}$.\\
		Ta có $\mathrm{P}(B) = \mathrm{P}\left(A \right) \cdot \mathrm{P}\left(B \mid A \right) + \mathrm{P}\left(\overline{A} \right) \cdot \mathrm{P}\left(B \mid \overline{A} \right) = \dfrac{2}{3} \cdot  \dfrac{55}{91} + \dfrac{1}{3} \cdot \dfrac{90}{91} = \dfrac{200}{273}$.\\
		Lại có $\mathrm{P}\left(A \mid B \right) = \dfrac{\mathrm{P}\left(A \right) \cdot P  \left(B \mid A \right)}{\mathrm{P}\left(B \right)} = \dfrac{9}{14} = 0{,}55$.
	} 
\end{bt}

% nguồn: Lop_12_Dai_so_Chuong_6
\begin{bt}%[ma:8394b6fa]
	%[2D6V1-4] 
	Lớp 12A có $37$ học sinh, trong đó có $15$ học sinh thích môn Tin học, $20$ học sinh thích môn Tiếng Anh, $10$ học sinh không thích môn nào trong hai môn trên. Chọn ngẫu nhiên $1$ học sinh. Xác suất chọn được học sinh thích môn Tin học, biết học sinh đó  thích môn Tiếng Anh, là bao nhiêu?
	 \shortans{0,4}
	\loigiai{ 
		Xét  các biến cố: $A\colon$\lq \lq Chọn được học sinh thích môn Tin học\rq \rq;\\
		$B\colon$\lq \lq Chọn được học sinh thích môn Tiếng  Anh\rq \rq.\\
		Khi  đó, $\mathrm{P}(A)=\dfrac{15}{37}$; $\mathrm{P}(B)=\dfrac{20}{37}$; $\mathrm{P}(A \cup B)=1-\dfrac{10}{37}=\dfrac{27}{37}$.\\
		Suy ra \[\mathrm{P}(A \cap B)=\mathrm{P}(A)+\mathrm{P}(B)-\mathrm{P}(A \cup B)=\frac{15}{37}+\frac{20}{37}-\frac{27}{37}=\frac{8}{37}.\] 
		Vậy xác suất chọn được học sinh thích môn  Tin học, biết học sinh đó thích môn
		Tiếng Anh,  là \[\mathrm{P}(A \mid B)=\dfrac{\dfrac{8}{37}}{\dfrac{20}{37}}=0{,}4.\]
	} 
\end{bt}

% nguồn: Lop_12_Dai_so_Chuong_6
\begin{bt}%[ma:d1edbdfe]
	 %[2D6V2-2]
	Có hai thùng I và II chứa các sản phẩm có khối lượng và hình dạng như nhau. Thùng I có $5$ chính phẩm và $4$ phế phẩm, thùng $2$ có $6$ chính phẩm và $8$ phế phẩm. Lấy ngẫu nhiên $1$ sản phẩm từ thùng I sang thùng II. Sau đó, lấy ngẫu nhiên $1$ sản phẩm từ thùng II để sử dụng. Xác  suất lấy được chính phẩm từ thùng II là bao nhiêu (làm tròn kết quả đến hàng phần trăm)?
	 \shortans{0,44}
	\loigiai{ 
		Xét các biến cố $A$: \lq \lq Lấy được $1$ chính phẩm từ thùng I sang thùng II\rq \rq; ${B}$: \lq \lq Lấy được $1$ chính phẩm từ thùng II\rq  \rq.\\
		 Khi đó, \[\mathrm{P}(A)=\frac{5}{9}; \mathrm{P}\big(\overline{A}\big)=\frac{4}{9}; \mathrm{P}(B \mid A)=\frac{7}{15}; \mathrm{P}\big(B \mid \overline{A}\big)=\frac{6}{15}=\frac{2}{5}.\]
		Theo  công thức xác suất toàn phần, xác suất của biến cố $B$ là
		\[\mathrm{P}(B)=\mathrm{P}(A) \cdot \mathrm{P}(B \mid A)+\mathrm{P}\big(\overline{A}\big) \cdot \mathrm{P}\big(B \mid \overline{A}\big)=\frac{5}{9} \cdot \frac{7}{15}+\frac{4}{9} \cdot \frac{2}{5} \approx 0{,}44.\] 
	} 
\end{bt}

% nguồn: Lop_12_Dai_so_Chuong_6
\begin{bt}%[ma:8fcef852]
	%[2D6V2-2] 
	Có $40$ tấm thẻ kích thước như nhau và đánh số thứ tự lần lượt từ $1$ đến  $40$ (mỗi tấm thẻ chỉ ghi một số nguyên dương, hai thẻ khác nhau ghi hai số khác nhau). Một người lần lượt rút hai thẻ (rút không hoàn lại). Tính xác suất lần thứ hai rút được thẻ ghi số nguyên~tố.
	 \shortans{0,3}
	\loigiai{ 
		 Xét các biến cố:\\
		 $A \colon$ \lq\lq Lần thứ nhất rút được thẻ ghi số nguyên~tố\rq\rq;\\
		$A \colon$ \lq\lq  Lần thứ hai rút được thẻ ghi số nguyên~tố\rq\rq.\\
		Từ $1$ đến $40$ có $12$ số nguyên~tố nên $P(A) = \dfrac{12}{40} = 0{,}3$ suy ra $P\left( \overline{A}  \right) = 1 - 0{,}3 = 0{,}7$.\\
		 Vì rút không hoàn lại nên $P \left(B \big| A \right) = \dfrac{11}{39}$ và $P \left(B \big| \overline{A} \right) = \dfrac{12}{39} = \dfrac{4}{13}$.\\
		Theo công thức xác suất  toàn phần, ta có
		\begin{align*}
			P\left(B \right)  = P\left(A \right) \cdot P \left(B \big| A \right) + P\left(\overline{A} \right) \cdot P \left(B \big| \overline{A} \right) = 0{,}3 \cdot \dfrac{11}{39} + 0{,}7 \cdot \dfrac{4}{13} = 0{,}3.
		\end{align*}
	} 
\end{bt}

% nguồn: Lop_12_Dai_so_Chuong_6
\begin{bt}%[ma:e4c8ec03]
	%[2D6V2-2] 
	Có hai chiếc hộp, hộp $I$ có $10$ quả bóng màu đỏ và $4$ quả bóng màu vàng, hộp $II$ có $10$ quả bóng màu đỏ và $8$ quả bóng màu vàng, các quả bóng có cùng kích thước và khối lượng. Lấy ngẫu nhiên một quả bóng từ hộp $I$ bỏ vào hộp $II$. Sau đó, lấy ra ngẫu nhiên một quả bóng từ hộp $II$. Tính xác  suất để quả bóng được lấy ra từ hộp $II$ là quả bóng được chuyền từ hộp $I$ sang, biết rằng quả bóng đỏ có màu đỏ (làm tròn kết quả đến hàng phần trăm).
	 \shortans{0,07}
	 \loigiai{
		Gọi $A$ là biến cố \lq\lq Quả bóng  lấy ra từ hộp $II$ là quả bóng được chuyển từ hộp $I$ sang\rq\rq.\\
		Gọi $B$ là biến cố \lq\lq Quả bóng lấy ra từ hộp $II$ là quả bóng  màu đỏ\rq\rq.
		\begin{center}
			\resizebox{0.8\linewidth}{!}{\begin{tikzpicture}[line join = round,line cap = round,>=stealth]
		\begin{scope}[ethicknode/.style={draw,rounded corners=5pt}]
			\coordinate (A)  at (0,0);
			\draw (A) node[left,text width=4cm] {Lấy một quả từ hộp $I$ bỏ  vào hộp $II$};
			\def \gocA{30} 
			\def  \kcA{4}
			\node  (B1) at ($(A)+(\gocA:\kcA)$){Vàng};
			\node (B2)  at ($(A)+(-\gocA:\kcA)$){Đỏ};
			 \def \gocB{0}
			 \def \kcB{4}
			\node (B11) at ($(B1)+(\gocB:\kcB)$){$AB$}; 
			\node (B21) at ($(B2)+(\gocB:\kcB)$){$AB$}; 
		\end{scope}
		\begin{scope}[ethicknode/.style={midway,sloped},ethickpath/.style={->}]
			\draw (A)--(B1)  node[above]{$\dfrac{2}{7}$};
			\draw  (A)--(B2) node[below]{$\dfrac{5}{7}$};
			\draw (B1)--(B11)  node[above]{$0$};
			\draw  (B2)--(B21) node[above]{$\dfrac{1}{19}$};
		\end{scope}
	\end{tikzpicture}
			} 
		\end{center}
		 Ta có
		$\mathrm{P}(AB)=\dfrac{5}{7} \cdot  \dfrac{1}{19}=\dfrac{5}{133}$.
		\begin{center}
			\resizebox{0.8\linewidth}{!}{\begin{tikzpicture}[line join = round,line cap = round,>=stealth,scale=1]
		\begin{scope}[ethicknode/.style={draw,rounded corners=5pt}]
			\coordinate (A) at  (0,0);
			\draw (A) node[text width=4cm,anchor=east] {Lấy một quả từ hộp $I$ bỏ vào hộp $II$}; 
			\def  \gocA{30}
			\def  \kcA{4}
			\node (B1)  at ($(A)+(\gocA:\kcA)$){Vàng};
			\node (B2)  at ($(A)+(-\gocA:\kcA)$){Đỏ};
			 \def \gocB{0}
			\def \kcB{4} 
			\node (B11) at ($(B1)+(\gocB:\kcB)$){$B$}; 
			\node (B21)  at ($(B2)+(\gocB:\kcB)$){$B$};
		\end{scope}
		\begin{scope}[ethicknode/.style={midway,sloped},ethickpath/.style={->}]
			\draw  (A)--(B1) node[above]{$\dfrac{2}{7}$};
			\draw  (A)--(B2) node[below]{$\dfrac{5}{7}$};
			\draw (B1)--(B11)  node[above]{$\dfrac{10}{19}$};
			\draw (B2)--(B21) node[above]{$\dfrac{11}{19}$}; 
		\end{scope}
	\end{tikzpicture}
			} 
		\end{center}
		$\mathrm{P}(B)=\dfrac{2}{7}  \cdot \dfrac{10}{19}+\dfrac{5}{7} \cdot \dfrac{11}{19}=\dfrac{75}{133}$.\\
		Vậy $\mathrm{P}(A  \mid B)=\dfrac{\mathrm{P}(AB)}{\mathrm{P}(B)}=\dfrac{1}{15} \approx 0{,}07$.
	} 
\end{bt}

% nguồn: Lop_12_Dai_so_Chuong_6
\begin{bt}%[ma:28bd4e9d]
	%[2D6V1-4] 
	Hai người ngang tài ngang sức tranh chức vô địch của cuộc thi cờ tướng. Người giành chiến thắng là người đầu tiên thắng được 5 ván cờ. Tài thời điểm người chơi thứ nhất đã thắng 4 ván và người chơi thứ hai mới thắng 2 ván, xác suất để người chơi thứ nhất giành chiến thắng là $P=\dfrac{a}{b},$ $(a,b\in \mathbb{N}^*;$ $a,b$ nguyên~tố cùng  nhau). Tính $a^2+b$?
	\shortans{57} 
	 \loigiai{
		Hai người cần chơi tối  đa $9$ ván cờ. Do đã chơi được $6$ ván và người thứ nhất đã thắng $4$ ván nên người thứ nhất giành chiến thắng chung cuộc khi trong ba ván còn lại thắng 1 ván.\\
		Gọi $A _i$ là biến cố: Người thứ nhất thắng ván thứ $ i$ trong ba ván  còn lại, $ i=1{,}2,3$.\\
		Các biến  cố $A_1,A_2,A_3$ độc lập và $P(A_i)=P(\bar{A}_i)=\dfrac{1}{2},i=1{,}2,3 $.\\
		Gọi $A$ là biến cố: Người  thứ nhất giành chiến thắng chung cuộc\\
		Khi  đó $A=A_1\cup \bar{A}_1A_2\cup \bar{A}_1\bar{A}_2A_3$\\
		Áp dụng công thức cộng và công thức nhân xác suất, ta có\\ 
		$$\begin{aligned}
		 &P(A)=P(A_1)+P(\bar{A}_1{A_2})+P(\bar{A}_1\bar{A}_2A_3)\\ &=P(A_1)+P(\bar{A}_1)P({A_2})+P(\bar{A}_1)P(\bar{A}_2)P(A_3)\\ &=\dfrac{1}2+\dfrac{1}{2^2}+\dfrac{1}{2^3}=\dfrac{7}8.
	\end{aligned}$$
Vậy  $ a=7,b=8\Rightarrow {a^2}+b=57$.}
\end{bt}

% nguồn: Lop_12_Dai_so_Chuong_6
\begin{bt}%[ma:ef0ba636]
	 %[2D6V1-2]
	Có hai chiếc hộp, hộp $I$ có $6$ bi đỏ và $4$ bi trắng, hộp $II$ có  $7$ bi đỏ và $3$ bi trắng, các bi có cùng kích thước và khối lượng. Lấy ngẫu nhiên từ mỗi hộp ra hai viên bi. Tính xác suất để lấy được ít nhất một bi đỏ từ hộp $I$, biết rằng trong bốn bi lấy ra số bi đỏ bằng số bi trắng.
	\shortans{0,813}%CẦN KIỂM  TRA ĐÁP ÁN
	 \loigiai{
		Gọi $A$ là biến cố lấy được ít nhất một bi đỏ từ hộp $I$.\\ 
		$B$ là  biến cố bốn bi lấy ra có số bi đỏ bằng số bi trắng.\\
		Ta có xác suất cần tìm là $P\left(A \mid B \right)  = \dfrac{P(AB)}{P(B)}$.\\
		Với $P(B) = \dfrac{n(B)}{n(\Omega)}= \dfrac{\mathrm{C}^{2}_{6} \cdot \mathrm{C}^{2}_{3} + \mathrm{C}^{2}_{4} \cdot \mathrm{C}^{2}_{7} + \left(6 \cdot 4 \right) \cdot \left(7 \cdot 3 \right)}{\mathrm{C}^{2}_{10} \cdot \mathrm{C}^{2}_{10}} = \dfrac{1}{3}$.\\ 
		Lại có $P(AB) = \dfrac{\left(6 \cdot 4 \right) \cdot \left(7 \cdot 3 \right)}{\mathrm{C}^{2}_{10} \cdot \mathrm{C}^{2}_{10}} + \dfrac{\mathrm{C}^{2}_{6} \cdot \mathrm{C}^{2}_{3}}{\mathrm{C}^{2}_{10} \cdot  \mathrm{C}^{2}_{10}} = \dfrac{61}{225}$.\\
		Khi đó xác suất cần tìm  là $P\left(A \mid B \right) = \dfrac{P(AB)}{P(B)} = \dfrac{61}{75} \approx 0{,}813$.
	} 
\end{bt}

% nguồn: Lop_12_Dai_so_Chuong_6
\begin{bt}%[ma:6968eb4c]
	 %[2D6V2-3]
	Trong một kỳ thi tốt nghiệp trung học phổ thông, một tỉnh $X$ có $80\%$ học sinh lựa chọn  tổ hợp $A00$ (gồm các môn Toán, Vật lý, Hóa học).
	Biết rằng: Nếu học sinh chọn tổ hợp $A00$, xác suất học sinh đó đỗ đại học là $0{,}6$. Nếu học sinh không chọn tổ hợp $A00$, xác suất học sinh đó đỗ đại học là $0{,}7$. Chọn ngẫu nhiên một học sinh của  tỉnh $X$ đã tốt nghiệp trung học phổ thông trong kỳ thi trên.
	Biết rằng học sinh đã đỗ đại học, tính xác suất để học sinh đó chọn tổ hợp $A00$  (làm tròn đến hàng phần trăm).
	\shortans{0,77} 
	 \loigiai{
		Gọi $A$ là biến cố: \lq\lq Học sinh chọn tổ hợp $A00$ \rq\rq, và $B$ là biến cố: \lq\lq Học  sinh đó đỗ đại học \rq\rq.\\
		Theo đề bài,  ta có các giá trị xác suất:
		 \[
		\mathrm{P}(A) = 0{,}8, \quad \mathrm{P}(\overline{A}) = 0{,}2, \quad \mathrm{P}(B|A) = 0{,}6, \quad \mathrm{P}(B|\overline{A}) = 0{,}7. 
		\] 
		Tính  xác suất $\mathrm{P}(B)$
		Sử dụng công thức xác  suất toàn phần:
		 \[
		\mathrm{P}(B)  = \mathrm{P}(A)\mathrm{P}(B|A) + \mathrm{P}(\overline{A})\mathrm{P}(B|\overline{A}).
		 \]
		Thay các giá  trị vào:
		\[ 
		\mathrm{P}(B) = 0{,}8  \cdot 0{,}6 + 0{,}2 \cdot 0{,}7 = 0{,}62.
		 \]
		 Tính xác suất $\mathrm{P}(A|B)$.\\
		Sử dụng định lý  Bayes:
		\[ 
		\mathrm{P}(A|B) = \dfrac{\mathrm{P}(A)\mathrm{P}(B|A)}{\mathrm{P}(B)}. 
		 \]
		 Thay các giá trị vào:
		\[ 
		\mathrm{P}(A|B) =  \dfrac{0{,}8 \cdot 0{,}6}{0{,}62} = \dfrac{24}{31}.
		\] 
		Xác suất học sinh chọn tổ hợp $A00$ khi đã đỗ đại  học là
		\[ 
		\mathrm{P}(A|B) =  0{,}77 \quad \text{ (làm tròn đến 2 chữ số thập phân)}.
		\] 
		 Có thể sử dụng sơ đồ cây.
	} 
\end{bt}

% nguồn: Lop_12_Dai_so_Chuong_6
\begin{bt}%[ma:3b01daa7]
	 %[2D6V2-2]
	Tỉ lệ bị bệnh cúm tại một địa phương bằng $0{,}25$. Khi thực hiện xét nghiệm chẩn đoán, nếu người có bệnh cúm thì khả năng phản ứng dương tính là $96 \%$, nếu người không bị bệnh cúm thì khả năng phản ứng dương tính $8 \%$. Chọn ngẫu nhiên  $1$ người tại địa phương đó. Xác suất người được chọn có phản ứng dương tính là bao nhiêu?
	\shortans{0,3} 
	 \loigiai{
		Xét các biến cố $A$: \lq  \lq Chọn được người bị bệnh cúm\rq \rq;
		$B$: \lq \lq Chọn được người có phản ứng dương  tính\rq \rq.\\
		Khi đó, $\mathrm{P}(A)=0{,}25$; $\mathrm{P}\big(\overline{A}\big)=0{,}75$; $\mathrm{P}(B  \mid A)=0{,}96$; $\mathrm{P}\big(B \mid \overline{A}\big)=0{,}08$.\\
		Theo công thức xác suất toàn phần, xác suất của biến  cố $B$ là
		\[ 
		\mathrm{P}(B)=\mathrm{P}(A) \cdot \mathrm{P}(B \mid A)+\mathrm{P}\big(\overline{A}\big)  \cdot \mathrm{P}(B \mid \overline {A})=0{,}25 \cdot 0{,}96+0{,}75 \cdot 0{,}08=0{,}3.
		 \]
	} 
\end{bt}

% nguồn: Lop_12_Dai_so_Chuong_6
\begin{bt}%[ma:31925b12]
	 %[2D6V2-2]
	Cho hai chuồng thỏ, chuồng thứ nhất có $5$ con thỏ đen và $10$ con thỏ trắng, chuồng thứ hai có $3$ con thỏ trắng và $7$ con thỏ đen, các con  thỏ có cùng khối lượng. Bắt ngẫu nhiên một con thỏ từ chuồng thứ nhất cho vào chuồng thứ hai. Sau đó bắt ngẫu nhiên một con thỏ từ chuồng thứ hai thì được một con thỏ trắng. Tính xác suất để thỏ trắng này là của chuồng thứ nhất mang sang (làm tròn kết quả đến hàng phần trăm).
	 \shortans{0,18}
	 \loigiai{
		Gọi $\Omega $ là không gian mẫu, ta có $n(\Omega)=15\cdot 10=150$.\\ 
		Gọi $A$ là biến cố: \lq\lq Bắt được con  thỏ ở chuồng thứ hai mà con thỏ đó được chuyển sang từ chuồng thứ nhất\rq\rq;\\
		$B$ là biến cố: \lq\lq Bắt được thỏ từ chuồng thứ  hai là thỏ màu trắng\rq\rq.\\
		Ta có $\mathrm{P}(B)=\dfrac {10\cdot 4+5\cdot 3}{150}=\dfrac  {55}{150}$.\\
		Biến cố $AB$: \lq\lq Con thỏ bắt được từ chuồng thứ hai có màu trắng được chuyển từ chuồng  thứ nhất sang\rq\rq.\\
		$\mathrm{P}(AB)=\dfrac  {10\cdot 1}{150}=\dfrac {10}{150}$.\\
		Xác suất để bắt  được thỏ trắng ở chuồng thứ hai được mang từ chuồng thứ nhất sang là\\
		$\mathrm{P}\left( A\mid B \right)=\dfrac {\mathrm{P}\left( AB \right)}{\mathrm{P}(B)}=\dfrac {\dfrac {10}{150}}{\dfrac {55}{150}}=\dfrac  {2}{11}\approx 0{,}18$.\\
	 }
\end{bt}

% nguồn: Lop_12_Dai_so_Chuong_6
\begin{bt}%[ma:5526820a]
	 %[2D6V1-4]
	Câu lạc bộ văn nghệ của trường Giải Phóng có $40$ bạn đều biết chơi ít nhất một trong hai loại đàn là organ và guitar, trong đó có $27$ bạn biết chơi đàn organ, $25$ bạn biết chơi đàn guitar. Chọn ngẫu nhiên $1$ bạn. Xác suất chọn được bạn biết chơi đàn organ, biết  bạn đó chơi được đàn guitar, là bao nhiêu?
	 \shortans{0,48}
	 \loigiai{
		Xét các biến cố sau: biến  cố ${A}\colon$\lq \lq Chọn được bạn biết chơi đàn organ\rq \rq;\\
		Biến  cố ${B}\colon$\lq \lq Chọn được bạn biết chơi đàn guitar\rq \rq.\\
		Khi đó,  $\mathrm{P}(A)=\dfrac{27}{40}=0{,}675$; $ \mathrm{P}(B)=\dfrac{25}{40}=0{,}625$; $\mathrm{P}(A \cup B)=1$.\\
		Suy ra  \[\mathrm{P}(A \cap B)=\mathrm{P}(A)+\mathrm{P}(B)-\mathrm{P}(A \cup B)=0{,}675+0{,}625-1=0{,}3.\]
	Vậy xác suất chọn được bạn biết chơi đàn organ, biết bạn đó chơi được đàn guitar,  là \[\mathrm{P}(A \mid B)=\dfrac{\mathrm{P}\big(A\cap B\big)}{\mathrm{P}(B)}=\dfrac{0{,}3}{0{,}625}=0{,}48.\]}
\end{bt}

% nguồn: Lop_12_Dai_so_Chuong_6
\begin{bt}%[ma:724be0cf]
	 %[2D6H1-2]
	Một khu dân cư có $60 \%$ các hộ gia đình có không quá $4$ thành viên. Trong các gia đình có không quá $4$ thành viên, có $20 \%$ gia đình có ba thế hệ cùng chung sống; trong các gia đình có trên $4$ thành viên, có $70  \%$ gia đình có ba thế hệ cùng chung sống. Chọn ngẫu nhiên $1$ hộ gia đình trong khu dân cư. Biết rằng gia đình đó có ba thế hệ cùng chung sống, tính xác suất để gia đình đó có trên $4$ thành viên.
	\shortans{0,4} 
	\loigiai{ 
		Gọi $M$ là biến cố \lq\lq Gia đình có trên $4$ thành viên\rq\rq, $N$ là biến cố \lq\lq Gia đình có $3$ thế hệ chung sống\rq\rq.\\ 
		Ta có $\mathrm P(M)=0{,}4$; $\mathrm P\left(\overline{M}\right)=0{,}6$; $\mathrm P\left(N \mid M\right)=0{,}7$; $\mathrm  P\left(N \mid \overline{M}\right)=0{,}2$.\\
		Theo công thức xác suất toàn phần,  ta có
		\[ 
		\mathrm P(N)=\mathrm P(M) \cdot \mathrm P\left(N \mid M\right)+\mathrm P\left(\overline{M}\right)  \cdot \mathrm P\left(N \mid \overline{M}\right)=0{,}4 \cdot 0{,}7+0{,}6 \cdot 0{,}2=0{,}4.
		 \]
		Vậy $\mathrm P\left(M \mid N\right)=\dfrac{\mathrm P(M) \cdot \mathrm P\left(N \mid  M\right)}{\mathrm P(N)}=\dfrac{0{,}4 \cdot 0{,}7}{0{,}4}=0{,}7$.
	 }
\end{bt}

% nguồn: Lop_12_Dai_so_Chuong_6
\begin{bt}%[ma:0e16d4f4]
	 %[2D6V1-4]
	Hai bạn Hải và Bình cùng tham dự một kì thi trắc nghiệm, vòng $1$ thi Toán, vòng $2$ thi Tiếng Anh. Mỗi vòng thi có $8$ mã đề được đánh số từ $1$ đến $8$. Mỗi bạn phải bốc thăm ngẫu nhiên $1$ đề Toán và $1$ đề Tiếng Anh. Xét biến  cố ${A}\colon$\lq \lq Hai bạn có chung mã đề ở duy nhất một vòng thi\rq \rq. Xác suất của biến cố $A$ là $\dfrac{a}{b}$ với $a$, $b$ là các số tự nhiên khác $0$, $b<50$. Giá trị của $a+b$ là bao nhiêu?
	\shortans{39} 
	\loigiai{ 
		Số cách chọn đề ngẫu nhiên của  bạn Hải là $8^2$.\\
		Sau đó bạn Bình chọn tiếp,  có số cách chọn bằng $8^2$.\\
		Suy ra  không gian mẫu có số phần tử là $8^4$.\\
		Xét biến cố $A\colon$\lq  \lq Hai bạn có chung mã đề ở duy nhất một vòng thi\rq \rq.\\
		Giả sử hai bạn chung mã đề ở vòng $1$, khác mã đề ở vòng $2$ thì số cách bốc thăm là $8  \mathrm{A}_8^2$. \\
		Nếu hai  bạn chung mã đề ở vòng $2$, khác mã đề ở vòng $1$ thì số cách bốc thăm cũng là $8 \mathrm{A}_8^2$.\\
		Khi đó, số kết  quả thuận lợi của biến cố $A$ là $16 \mathrm{A}_8^2$.\\
		Xác  suất của biến cố $A$ là \[\mathrm{P}(A)=\frac{16 \mathrm{A}_8^2}{8^4}=\frac{7}{32}.\]
		 Suy ra $a=7$; $b=32$.\\
		Vậy  $a+b=39$.
	 }
\end{bt}

% nguồn: Lop_12_Dai_so_Chuong_6
\begin{bt}%[ma:244a8223]
	 %[2D6C2-4]
	Có hai chiếc hộp, hộp I có $5$ quả bóng màu trắng và $7$ quả bóng màu đỏ, hộp II có $10$ quả bóng màu trắng và $15$ quả bóng màu đỏ, các quả bóng có cùng kích thước và khối  lượng. Lấy ngẫu nhiên hai quả bóng từ hộp I bỏ vào hộp II. Sau đó, lấy ra ngẫu nhiên một quả bóng từ hộp II. Xác suất để quả bóng được lấy ra từ hộp II là quả bóng được chuyển từ hộp I sang, biết rằng quả bóng đó có màu trắng là $\dfrac{a}{b}$ (làm tròn tối giản). Tính $a + b$.
	\shortans{14} 
	\loigiai{ 
		Gọi $A$ là \lq\lq sự kiện quả bóng lấy ra từ hộp II là một trong $2$ quả bóng được chuyển từ hộp  I\rq\rq.\\
		$B$ là \lq\lq sự kiện quả bóng lấy  ra từ hộp II có màu trắng\rq\rq.\\
		 Khi đó cần tính $$\mathrm{P}(A \mid B)=\dfrac{\mathrm{P}(A\cap B)}{\mathrm{P}(B)}.$$
		Tính $\mathrm{P}(A \cap B)$. 
		\begin{eqnarray*}
			\mathrm{P}(A \cap B)&=&P(2 \text {quả trắng}) \cdot 1+P(1 \text {quả trắng, } 1 \text {quả đỏ}) \cdot \dfrac{1}{2}+P(2 \text  {quả đỏ}) \cdot 0 \\
			&=&\dfrac{10}{66} \cdot 1+\dfrac{35}{66} \cdot \dfrac{1}{2}+\dfrac{21}{66} \cdot 0 \\ 
			 &=&\dfrac{55}{132}.
		\end{eqnarray*}
		Tính $\mathrm{P}(B)$, 
		gọi \\$A_1$: Lấp được $2$ bóng  trắng.\\
		$A_2$: Lấp được $1$ bóng đỏ,  $1$ bóng trắng.\\
		$A_3$: Lấp được $2$ bóng  đỏ.\\
		Ta có 
		\begin{eqnarray*}
			\mathrm{P}(B)&=&\mathrm{P}(A_1)\cdot  \mathrm{P}(B|A_1)+\mathrm{P}(A_2)\cdot \mathrm{P}(B|A_2)+\mathrm{P}(A_3)\cdot \mathrm{P}(B|A_3)\\
			&=&\dfrac{\mathrm{C}_5^2}{\mathrm{C}_{12}^2}\cdot \dfrac{12}{27}+\dfrac{\mathrm{C}_5^1\cdot \mathrm{C}_7^1}{\mathrm{C}_{12}^2}\cdot \dfrac{11}{27}+\dfrac{\mathrm{C}_7^2}{\mathrm{C}_{12}^2}\cdot \dfrac{10}{27}\\ 
			 &=&\dfrac{65}{162}.
		\end{eqnarray*}
		Ta được $$\mathrm{P}(A \mid B)=\dfrac{\mathrm{P}(A\cap  B)}{\mathrm{P}(B)}=\dfrac{\dfrac{55}{132}}{\dfrac{65}{162}}=\dfrac{1}{13}.$$
		Vậy  $a=1$, $b=13$ suy ra $a+b=14$.\\
		\textbf{Nhận xét:} Bài này  giải cách này nhanh gọn hơn dùng Bayes.
	} 
\end{bt}

% nguồn: Lop_12_Dai_so_Chuong_6
\begin{bt}%[ma:e0c02043]
	 %[2D6V2-3]
	Thực hiện khảo sát tại một địa phương mà số trẻ em  nam gấp $1{,}5$ lần số trẻ em nữ, có $8 \%$ số trẻ em nam bị hen phế quản, $5 \%$ số trẻ em nữ bị hen phế quản. Chọn ngẫu nhiên $1$ trẻ em. Giả sử trẻ em được chọn bị hen phế quản. Xác suất chọn được trẻ em nam là bao nhiêu (làm tròn kết quả đến hàng phần mười)?
	\shortans{0,7} 
	\loigiai{ 
		Xét các biến  cố $A\colon$\lq \lq Chọn được trẻ em nam\rq \rq;
		 $B\colon$\lq \lq Chọn được trẻ em bị hen phế quản\rq \rq.\\
		Khi  đó, $\mathrm{P}(A)=\frac{1{,}5}{1+1{,}5}=0{,}6$; $\mathrm{P}\big(\overline{A}\big)=0{,}4$; $\mathrm{P}(B \mid A)=0{,}08$; $\mathrm{P}\big(B \mid \overline{A}\big)=0{,}05$.\\
		Theo công thức xác suất toàn phần,  ta có
		 \[
		\mathrm{P}(B)=\mathrm{P}(A) \cdot  \mathrm{P}(B \mid A)+\mathrm{P}\big(\overline{A}\big) \cdot \mathrm{P}\big(B \mid \overline{A}\big)=0{,}6 \cdot 0{,}08+0{,}4 \cdot 0{,}05=0{,}068.
		 \]
		Theo công thức Bayes, xác suất chọn được trẻ em nam, biết trẻ em đó  bị hen phế quản là \[\mathrm{P}(A \mid B)=\frac{\mathrm{P}(A) \cdot \mathrm{P}(B \mid A)}{\mathrm{P}(B)}=\frac{0{,}6 \cdot 0{,}08}{0{,}068} \approx 0{,}7.\]
	} 
\end{bt}

% nguồn: Lop_12_Dai_so_Chuong_6
\begin{bt}%[ma:070fe173]
	%[2D6H1-2] 
	Năm bạn $A$, $B$, $C$, $D$, $E$ xếp thành một hàng ngang theo thứ tự ngẫu nhiên. Tính xác suất để $A$ đứng cạnh $C$, biết rằng  $A$ không đứng cạnh $E$.
	\shortans{0,5} 
	 \loigiai{
		Xếp $5$ bạn thành hàng ngang nên $n(\Omega)=5!=120$ cách.\\ 
		Gọi $Y$ là biến cố $A$ không đứng cạnh $E$, $X$ là biến cố $A$ đứng cạnh $C$. Ta cần tính $\mathrm{P}(X|Y)$.\\ 
		Số cách xếp $A$ đứng cạnh $E$ là  $n\left(\overline Y\right)=4!\cdot2!=48$. Suy ra
		$$\mathrm{P}(Y)=1-\mathrm{P}\left(\overline Y\right)=1-\dfrac{48}{120}=\dfrac{3}{5}.$$ 
		Xét $XY$ là \lq\lq Xếp $A$ đứng cạnh $C$ nhưng $A$ không  cạnh $E$\rq\rq.\\
		Do $A$ cạnh $C$ nên ta xem $AC$ là 1 vị trí cần xếp trong hàng, khi đó chỉ cần xếp $4$  vị trí. Giả sử các vị trí cần xếp là $a$, $b$, $c$, $d$.
		\begin{itemize}
			\item \textbf{Trường hợp 1.} Xếp $AC$ vào hàng. 
			\begin{itemize}
				\item Nếu $AC$ ở  vị trí $a$, ta có $3$ cách chọn xếp vị trí cho $E$, và có 2 cách xếp $B$, $D$ vào 2 chỗ còn lại. Khi đó có $3\cdot 2=6$ cách.
				\item Nếu $AC$ ở vị trí $b$, $c$, hoặc $d$, ta có $2$ cách chọn xếp vị trí cho $E$, và có 2 cách  xếp $B$, $D$ vào 2 chỗ còn lại. Khi đó có $3\cdot 2\cdot2=12$ cách.
			\end{itemize}
			 Trường hợp 1 có $6+12=18$ cách.
			\item \textbf{Trường hợp 2.}  Xếp $CA$ vào hàng. Lập luận tương tự ta cũng có $18$ cách xếp.
		\end{itemize}
		Suy ra $n(XY)=18+18=36$, do đó  $\mathrm{P}(XY)=\dfrac{36}{120}=\dfrac{3}{10}$.\\
		 Vậy $\mathrm{P}(X|Y)=\dfrac{\mathrm{P}(XY)}{\mathrm{P}(Y)}=\dfrac{\tfrac{3}{10}}{\tfrac{3}{5}}=\dfrac{1}{2}=0{,}5$.
	 }
\end{bt}
